\documentclass[11pt,a4paper]{amsart}
\usepackage[T1]{fontenc}
\usepackage[utf8]{inputenc}
\usepackage{lmodern}
\usepackage[american]{babel}
\usepackage{amsmath,amssymb,amsthm,mathtools}
\usepackage[all,cmtip]{xy}
\usepackage{microtype}
\usepackage{enumitem}
\usepackage{url}
\usepackage{xspace}
\usepackage{geometry}
\geometry{margin=1.1in}

% --- SGA 1 macros (subset used in Exposé VI) ---
\DeclareMathOperator{\Aut}{Aut}
\DeclareMathOperator{\SheafAut}{\mathbf{Aut}}
\newcommand{\Cat}{\mathbf{Cat}}
\newcommand{\Ens}{\mathbf{Ens}}
\newcommand{\cart}{\mathrm{cart}}
\newcommand{\Fib}{\mathrm{Fib}}
\DeclareMathOperator{\Fl}{\mathrm{Fl}}
\DeclareMathOperator{\Hom}{Hom}
\DeclareMathOperator{\SheafHom}{\mathbf{Hom}}
\newcommand{\id}{\mathrm{id}}
\DeclareMathOperator{\Isom}{Isom}
\DeclareMathOperator*{\Lim}{Lim}
\DeclareMathOperator{\Ob}{\mathrm{Ob}}
\newcommand{\Sch}{\mathbf{Sch}}
\newcommand{\Univ}{\mathbf{Univ}}
\DeclareMathOperator{\CatFl}{\mathbf{Fl}}
\DeclareMathOperator{\Ker}{Ker}
\DeclareMathOperator{\pr}{pr}
\renewcommand{\H}{\mathrm{H}}
\newcommand{\isomto}{\overset{\sim}{\longrightarrow}}
\newcommand{\isomfrom}{\overset{\sim}{\longleftarrow}}
\let\cal\mathcal
\let\tilde\widetilde

\def\to{\mathchoice{\longrightarrow}{\rightarrow}{\rightarrow}{\rightarrow}}
\def\from{\mathchoice{\longleftarrow}{\leftarrow}{\leftarrow}{\leftarrow}}
\def\mto{\mathchoice{\longmapsto}{\mapsto}{\mapsto}{\mapsto}}
\def\To{\mathchoice{\Longrightarrow}{\Rightarrow}{\Rightarrow}{\Rightarrow}}
\def\lto#1{\mathchoice{\xrightarrow{\scriptstyle~#1~}}{\stackrel{#1}{\longrightarrow}}{}{}}
\def\lfrom#1{\mathchoice{\xleftarrow{\scriptstyle~#1~}}{\stackrel{#1}{\longleftarrow}}{}{}}
\def\quoi{}
\def\cf{{cf.}\kern.3em}
\def\Cf{{Cf.}\kern.3em}
\def\ie{{i.e.}\kern.3em}
\def\iev{{i.e.},\xspace}
\def\resp{resp.\xspace}
\def\loccit{loc.\kern.2pt cit.\xspace}
\def\ptbl{.\kern.2em }
\newcommand{\nno}{no.~}
% Kernel \Ref is a tagged-PDF command (star argument). Use \ref.
\def\dpl{\displaystyle}
\let\varprojLim\varprojlim

% theorems numbered as in SGA: shared counter per section
\theoremstyle{plain}
\newtheorem{theorem}{Theorem}[section]
\newtheorem{proposition}[theorem]{Proposition}
\newtheorem{lemma}[theorem]{Lemma}
\newtheorem{corollary}[theorem]{Corollary}
\theoremstyle{definition}
\newtheorem{definition}[theorem]{Definition}
\theoremstyle{remark}
\newtheorem*{remarks}{Remarks}
\newtheorem*{remark}{Remark}

\numberwithin{equation}{section}
\setcounter{section}{-1}

\title[SGA~1, Expos\'e VI (English)]{Fibered categories and descent\\
{\normalsize SGA~1, Expos\'e VI}}
\author{A.\ Grothendieck}
\date{}
\thanks{English translation of Expos\'e~VI of
\textit{Rev\^etements \'etales et groupe fondamental} (SGA~1),
S\'eminaire de g\'eom\'etrie alg\'ebrique du Bois Marie, 1960--61,
recomposed SMF edition, arXiv:math/0206203v2.
Only this expos\'e is translated.}

\begin{document}
\begin{abstract}
This is an English translation of Expos\'e~VI,
``Cat\'egories fibr\'ees et descente'', of SGA~1
(A.\ Grothendieck, \textit{Rev\^etements \'etales et groupe fondamental},
1960--61). The mathematical content follows the slightly corrected SMF
recomposition (arXiv:math/0206203v2). Numbering of statements is that
of the original expos\'e.
\end{abstract}
\maketitle
\tableofcontents
\newpage

% ---- en-chunk1.tex ----
\label{VI}
\section{Introduction}

Contrary to what had been announced in
the introduction to the preceding expos\'e, it has
turned out to be impossible to do descent in the category
of preschemes, even in special cases, without having
first developed with sufficient care the language of
descent in general categories.

The notion of ``descent'' provides the general framework for all
procedures of ``gluing'' of objects, and
consequently of ``gluing'' of categories. The most
classical case of gluing is relative to the datum of a
topological space $X$ and a covering of~$X$ by open sets $X_i$; if
one is given for every $i$ a fiber space (say) $E_i$ over
$X_i$, and for every pair $(i,j)$ an isomorphism $f_{ji}$ of
$E_i|X_{ij}$ onto~$E_j|X_{ij}$ (where one sets $X_{ij}=X_i\cap X_j$),
satisfying a well-known transitivity condition (which one
writes in abbreviated form $f_{kj}f_{ji}=f_{ki}$), one
knows that there exists a fiber space $E$ over~$X$, defined up to
isomorphism by the condition that one have isomorphisms
$f_i\colon E|X_i\isomto E_i$, satisfying the relations
$f_{ji}=f_jf_i^{-1}$ (with the usual abuse of writing). Let $X'$
be the sum space of the $X_i$, which is thus a fiber space over
$X$ (\ie equipped with a continuous map $X'\to X$). One can
interpret more concisely the datum of the $E_i$
as a fiber space $E'$ over~$X'$, and the datum of the $f_{ji}$
as an isomorphism between the two inverse images (by the two
canonical projections) $E''_1$ and $E''_2$ of~$E'$ on~$X''=X'\times_X
X'$, the gluing condition then being able to be written as an
identity between isomorphisms of fiber spaces $E'''_1$ and
$E'''_3$ over the triple fibered product $X'''=X'\times_X X'\times_X
X'$ (where $E'''_i$ denotes the inverse image of~$E'$ on~$X'''$
by the canonical projection of index $i$). The construction of~$E$,
starting from~$E'$ and from~$f$, is a typical case
% original p. 146
of a procedure of ``descent''. Moreover, starting from a fiber
space $E$ over~$X$, one says that $X$ is ``locally trivial'', of
fiber $F$, if there exists an open covering $(X_i)$ of~$X$ such that
the $E|X_i$ are isomorphic to $F\times X_i$, or what amounts
to the same, such that the inverse image $E'$ of~$E$ on~$X'=\coprod_i
X_i$ is isomorphic to $X'\times F$.

Thus, the notion of ``gluing'' of objects like that of
``localization'' of a property, are tied to
the study of certain types of ``changes of base'' $X'\to X$.
In Algebraic Geometry, many other types of change
of base, and notably the faithfully flat morphisms $X'\to X$,
must be regarded as corresponding to a
procedure of ``localization'' relative to
preschemes, or other objects, ``over''~$X$. This type of
localization is used just as much as simple
topological localization (which is moreover a special case of it). The same holds
(to a lesser extent) in Analytic Geometry.

Most of the proofs, reducing to
verifications, are omitted or merely sketched; if need
be we specify the less evident diagrams
which arise in a proof.

\section{Universes, categories, equivalence of categories}
\label{VI.1}

To avoid certain logical difficulties, we shall admit
here the notion of \emph{Universe}, which is a set ``large enough''
that one does not leave it by the usual operations of
set theory; an ``\emph{axiom of Universes}'' guarantees
that every object lies in a Universe. For details, see
a book in preparation by C\ptbl Chevalley and the speaker\footnote{%
The definitive authors are C\ptbl Chevalley and P\ptbl Gabriel. The book is to appear in the year 3000. In the meantime, \cf also SGA~4 I.}. Thus, the symbol~$\Ens$
designates, not the category of all sets (a notion
which has no meaning), but the category of sets which
lie in a given Universe (which we shall not specify
here in the notation). Likewise, $\Cat$
will designate the category of categories lying in
the Universe in question, the ``morphisms'' from an object $X$ of~$\Cat$
to another $Y$ being by definition the \emph{functors}
from~$X$ to $Y$.

If
% original p. 147
$\cal{C}$ is a category, we denote by $\Ob(\cal{C})$
\emph{the set of objects} of~$\cal{C}$, by~$\Fl(\cal{C})$
\emph{the set of arrows} of~$\cal{C}$ (or morphisms of
$\cal{C}$). We shall thus write $X\in\Ob(\cal{C})$, avoiding
the common abuse of notation $X\in\cal{C}$. If $\cal{C}$ and $\cal{C}'$
are two categories, a \emph{functor} from~$\cal{C}$ to
$\cal{C}'$ will always be what is commonly called a
\emph{covariant} functor from~$\cal{C}$ to $\cal{C}'$; its datum implies
that of the target category and the source
category, $\cal{C}$ and $\cal{C}'$. The functors from~$\cal{C}$ to
$\cal{C}'$ form a set, denoted $\Hom(\cal{C},\cal{C}')$, which
is the set of objects of a category denoted
$\SheafHom(\cal{C},\cal{C}')$.
By definition, a \emph{contravariant functor} from~$\cal{C}$ to
$\cal{C}'$ is a functor from the \emph{opposite category}
$\cal{C}^\circ$
of~$\cal{C}$ to~$\cal{C}'$.

We shall admit the notion of \emph{projective limit} and of
\emph{inductive limit} of a functor $F\colon \cal{I}\to \cal{C}$, and
in particular the most common special cases of this notion:
cartesian products and fibered products, dual notions of
direct sums and of amalgamated sums, and the usual
formal properties of these operations.

For example, in the category $\Cat$ introduced above, the
projective limits (relative to categories $\cal{I}$
lying in the chosen Universe) exist; the set of objects
(\resp the set of arrows) of the projective limit
category $\cal{C}$ of the $\cal{C}_i$, is obtained by taking the
projective limit of the sets of objects (\resp of the sets of arrows)
of the categories $\cal{C}_i$. The best-known case is that of the
product of a family of categories. We shall constantly use
in the sequel the fibered product of two categories over a
third.

For everything concerning categories and functors, while
awaiting the book in preparation already mentioned, see
\cite{VI.1} (which is necessarily very incomplete, even as regards
the generalities sketched in the
present no.).

Let us take this opportunity to make explicit the notion of equivalence of
categories,
which is not presented in a satisfactory way
in~\cite{VI.1}. A functor $F\colon\cal{C}\to \cal{C}'$ is said to be
\emph{faithful}
(\resp \emph{fully faithful})
if for every pair of objects $S$, $T$ of~$\cal{C}$, the map
$u\mto F(u)$ from~$\Hom(S,T)$ to $\Hom(F(S),F(T))$ is injective
(\resp bijective). One says that $F$ is an \emph{equivalence} of
categories if
% original p. 148
$F$ is fully faithful, and if moreover every object $S'$ of
$\cal{C}'$ is isomorphic to an object of the form $F(S)$. One shows
that this amounts to the same as saying that there exists a functor $G$ from
$\cal{C}'$ to $\cal{C}$ \emph{quasi-inverse of} $F$, \iev such that
$GF$ is isomorphic to $\id_{\cal{C}}$. When this is so, the
datum of a functor $G\colon\cal{C}'\to \cal{C}$ and of an
isomorphism $\varphi\colon GF\to \id_{\cal{C}'}$ is equivalent to
the datum, for every $S'\in\Ob(\cal{C}')$, of a pair $(S,u)$
formed of an object $S$ of~$\cal{C}$ and an isomorphism $u\colon
F(S)\to S'$, namely $(G(S),\varphi(S))$. (With these notations, there exists
a unique functor $\cal{C}'\to\cal{C}$ having the given
map $S\mto G(S)$ as object-map, and such that the map
$S\mto \varphi(S)$ is a homomorphism of functors $FG\to
\id_{\cal{C}'}$). Finally, if $G$ is a functor quasi-inverse of~$F$,
and if one chooses isomorphisms $\varphi\colon
FG\isomto\id_{\cal{C}'}$ and $\psi\colon GF\isomto\id_{\cal{C}}$, then
the two compatibility conditions on~$\varphi$, $\psi$
stated in \cite[I.1.2]{VI.1} are in fact equivalent to one
another, and for every chosen isomorphism $\varphi$, there exists
a unique isomorphism $\psi$ such that the said conditions are
satisfied.

\section{Categories over another}
\label{VI.2}

Let $\cal{E}$ be a category in~$\Univ$, it is thus an object
of~$\Cat$, and one can consider the category
$\Cat_{/\cal{E}}$
of ``objects of~$\Cat$ over~$\cal{E}$''. An object of this
category is thus a functor
$$
p\colon \cal{F}\to\cal{E}
$$
One also says that the category~$\cal{F}$, equipped with such a functor,
is a \emph{category over}~$\cal{E}$, or an
$\cal{E}$\emph{-category}. One will thus call an $\cal{E}$-functor
from a category $\cal{F}$ over~$\cal{E}$ to a category
$\cal{G}$ over~$\cal{E}$, a functor
$$
f\colon\cal{F}\to\cal{G}
$$
such that
$$
qf=p
$$
where $p$ and $q$ are the projection functors for $\cal{F}$
\resp $\cal{G}$. The set of $\cal{E}$-functors $f$ from~$\cal{F}$
to $\cal{G}$ is thus in one-to-one correspondence with the set
of arrows with source $\cal{F}$ and target $\cal{G}$
in~$\Cat_{/\cal{E}}$, without
% original p. 149
there being an identity there, however (since the datum of an
$f$ as above does not determine $\cal{F}$ and $\cal{G}$ as
categories over~$\cal{E}$); but of course, as in any
other category~$\cal{C}_{/S}$, one will commonly make the abuse of
language consisting in identifying the $\cal{E}$-functors (in the
sense made explicit above) with arrows in a
category~$\Cat_{/\cal{E}}$.

One will denote by
$$
\Hom_{\cal{E}}(\cal{F},\cal{G})
$$
the set of $\cal{E}$-functors from~$\cal{F}$ to~$\cal{G}$. Of
course, a composite of~$\cal{E}$-functors is an
$\cal{E}$-functor (the composition in question corresponding by
definition to the composition of arrows
in~$\Cat_{/\cal{E}}$).

Consider now two $\cal{E}$-functors
$$
f,g\colon \cal{F}\to\cal{G}
$$
and a homomorphism of functors:
$$
u\colon f\to g
$$
One says that $u$ is an $\cal{E}$-\emph{homomorphism} or a
``\emph{homomorphism of~$\cal{E}$-functors}'', if for every
$\xi\in\Ob(\cal{F})$, one has
$$
q(u(\xi))=\id_{p(\xi)} \quoi\text{,}
$$
in words: setting
$S=p(\xi)=qf(\xi)=qg(\xi)\in\Ob(\cal{E})$,
the morphism
$$
u(\xi)\colon f(\xi)\to g(\xi)
$$
in $\cal{G}$ is an $\id_S$-morphism. (In
general, for every morphism $\alpha\colon T \to S$
in~$\cal{E}$, and every category $\cal{G}$ over
$\cal{E}$, a morphism $v$ in $\cal{G}$ is called an
$\alpha$-\emph{morphism} if $q(v)=\alpha$, $q$ denoting the
projection functor $\cal{G}\to \cal{E}$). If one has a third
$\cal{E}$-functor $h\colon \cal{F}\to \cal{G}$ and an
$\cal{E}$-homomorphism $v\colon g\to h$, then $vu$ is likewise
an $\cal{E}$-homomorphism. Thus,
% original p. 150
\emph{the $\cal{E}$-functors from~$\cal{F}$ to~$\cal{G}$, and
the $\cal{E}$-homomorphisms of such, form a subcategory
of the category $\SheafHom(\cal{F},\cal{G})$ of all functors
from~$\cal{F}$ to~$\cal{G}$, which one will call the category of
$\cal{E}$-functors from~$\cal{F}$ to $\cal{G}$ and which one will denote}
$$
\SheafHom_{\cal{E}/-}(\cal{F},\cal{G})
$$
It is also the kernel subcategory of the pair of functors
$$
\xymatrix@C=.5cm{R,S\colon\SheafHom(\cal{F},\cal{G})
\ar@<2pt>[r]\ar@<-2pt>[r]&\SheafHom(\cal{F},\cal{E})}\quoi\text{,}
$$
where $R$ is the constant functor defined by the object $p$ of
$\SheafHom(\cal{F},\cal{E})$, and where $S$ is the functor $f\mto
q\circ f$ defined by $q\colon\cal{G}\to\cal{E}$.

To finish these generalities, it remains to define the
natural pairings among the categories
$\SheafHom_{\cal{E}/-}(\cal{F},\cal{G})$ by composition of
$\cal{E}$-functors. In other words, one wants to define a
``composition functor'':
\begin{equation*}
\tag{i}
\label{eq:VI.2.i} {\SheafHom_{\cal{E}/-}(\cal{F},\cal{G})\times
\SheafHom_{\cal{E}/-}(\cal{G},\cal{H})\to
\SheafHom_{\cal{E}/-}(\cal{F},\cal{H})}
\end{equation*}
when $\cal{F}$, $\cal{G}$, $\cal{H}$ are three categories
over~$\cal{E}$, in such a way that this functor induce, for the
objects, the composition map $(f,g)\mto gf$ of
$\cal{E}$-functors $f\colon\cal{F}\to\cal{G}$ and $g\colon \cal{G}\to
\cal{H}$. For this, let us recall that one defines a canonical
functor:
\begin{equation*}
\tag{ii}
\label{eq:VI.2.ii}
{\SheafHom(\cal{F},\cal{G})\times\SheafHom(\cal{G},\cal{H})\to
\SheafHom(\cal{F},\cal{H})}
\end{equation*}
which, for the objects, is nothing other than the composition map
$(f,g)\mto gf$ of functors, and which transforms an arrow
$(u,v)$,~where
$$
u\colon f\to f'\quoi,\quad v\colon g\to g'
$$
are arrows in $\SheafHom(\cal{F},\cal{G})$ \resp in
$\SheafHom(\cal{G},\cal{H})$, into the arrow
$$
v*u\colon gf\to g'f'
$$
defined by the relation:
% original p. 151
$$
v*u(\xi)=v(f'(\xi))\cdot g(u(\xi))=g'(u(\xi))\cdot v(f(\xi))
$$
It is well known that one thus indeed obtains a homomorphism from~$gf$
to~$g'f'$, and that (for $f,g$ and $u,v$ variable) one thus obtains
a functor~\eqref{eq:VI.2.ii}, \ie that one has
\begin{equation*}
\tag{I}
\label{eq:VI.2.I} {\id_g*\id_f=\id_{gf}\quoi,}
\end{equation*}
\begin{equation*}
\tag{II}
\label{eq:VI.2.II} {(v'*u')\circ(v*u)=(v'\circ v)*(u'\circ u)}
\end{equation*}
Let us also recall that one has an associativity formula for the
canonical pairings~\eqref{eq:VI.2.ii}, which is expressed on the one hand
by the associativity $(hg)f=h(gf)$ of composition of functors,
and on the other hand by the formula
\begin{equation*}
\tag{III}
\label{eq:VI.2.III} {(w*v)*u=w*(v*u)}
\end{equation*}
for the composition products of homomorphisms of functors (where
$u\colon f\to f'$ and $v\colon g\to g'$ are as above, and where one
assumes given moreover a homomorphism $w\colon h\to h'$ of
functors $h,h'\colon \cal{H}\to\cal{K}$). I now say that
\emph{when~$\cal{F}$, $\cal{G}$ are $\cal{E}$-categories,
the canonical composition functor}~\eqref{eq:VI.2.ii} \emph{induces a
functor}~\eqref{eq:VI.2.i}. As one already knows that the
composite of two $\cal{E}$-functors is an $\cal{E}$-functor,
this amounts to saying that \emph{when $u\colon f\to f'$ and
$v\colon g\to g'$ are homomorphisms of~$\cal{E}$-functors,
then $v*u\colon gf\to g'f'$ is likewise a homomorphism of
$\cal{E}$-functors.} This follows indeed trivially from the
definitions. As the pairings~\eqref{eq:VI.2.i} are
induced by the pairings~\eqref{eq:VI.2.ii}, they satisfy the
same associativity property, expressed also
in the formulae $(hg)f=h(gf)$ and $(w*v)*u=w*(v*u)$ for
$\cal{E}$-functors and $\cal{E}$-homomorphisms of
$\cal{E}$-functors.

To complete the formulary~\eqref{eq:VI.2.I},
\eqref{eq:VI.2.II}, \eqref{eq:VI.2.III}, let us also recall the formulae:
\begin{equation*}
\tag{IV}
\label{eq:VI.2.IV} {v*\id_{\cal{F}}=v\quad\text{and}\quad\id_{\cal{G}}*u=u}\quoi,
\end{equation*}
where
% original p. 152
for simplicity one writes $v*f$ or $u*g$ instead of~$v*u$, when
$u$ \resp $v$ is the identity automorphism of~$f$ \resp $g$.

From the definition of the pairings~\eqref{eq:VI.2.i} it follows
that $\SheafHom_{\cal{E}/-}(\cal{F},\cal{G})$ \emph{is a functor in
$\cal{F}$,~$\cal{G}$,
of the product category
${\Cat_{/\cal{E}}}^\circ\times\Cat_{/\cal{E}}$ to the
category~$\Cat$}. If indeed $g\colon \cal{G}\to\cal{G}_1$ is
an $\cal{E}$-functor, \ie an object of
$\SheafHom_{\cal{E}/-}(\cal{G},\cal{G}_1)$, then setting
in~\eqref{eq:VI.2.i} $\cal{H}=\cal{G}_1$, there corresponds to it a
functor
$$
g_*\colon\SheafHom_{\cal{E}/-}(\cal{F},\cal{G})\to
\SheafHom_{\cal{E}/-}(\cal{F},\cal{G}_1)
$$
One defines analogously, for an $\cal{E}$-functor
$f\colon \cal{F}_1\to \cal{F}$, a functor
$$
f^*:\SheafHom_{\cal{E}/-}(\cal{F},\cal{G})\to
\SheafHom_{\cal{E}/-}(\cal{F}_1,\cal{G})
$$
For brevity, one also denotes these functors by the symbols
$f\mto g\circ f$ \resp $g\mto g\circ f$ (which designate
only, in fact, the corresponding maps on the sets
of objects). It follows from the associativity property
indicated above that in this way, one indeed obtains as
announced a functor
$\Cat_{/\cal{E}}^\circ\times\Cat_{/\cal{E}}\to\Cat$.

\section{Change of base in categories over~$\cal{E}$}
\label{VI.3}

As in $\Cat$ the projective limits (relative to
categories $\cal{I}$ elements of~$\Univ$) exist, the same
is true in $\Cat_{/\cal{E}}$, in particular cartesian
products exist there, which are interpreted as fibered
products in~$\Cat$. In accordance with the general
notations, if $\cal{F}$ and $\cal{G}$ are categories
over~$\cal{E}$, one denotes by
$$
\cal{F}\times_{\cal{E}}\cal{G}
$$
their product in~$\Cat_{/\cal{E}}$, \ie their fibered product
over~$\cal{E}$ in~$\Cat$,
% original p. 153
as a category over~$\cal{E}$. Thus,
$\cal{F}\times_{\cal{E}} \cal{G}$ is equipped with two $\cal{E}$-functors
$\pr_1$ and $\pr_2$, which define, for every category
$\cal{H}$ over~$\cal{E}$, a bijection
$$
\Hom_{\cal{E}}(\cal{H},\cal{F}\times_{\cal{E}}\cal{G})\isomto
\Hom_{\cal{E}}(\cal{H},\cal{F})\times\Hom_{\cal{E}}(\cal{H},\cal{G}).
$$
This bijection arises moreover from an isomorphism of
categories
$$
\SheafHom_{\cal{E}/-}(\cal{H},\cal{F}\times_{\cal{E}}
\cal{G})\isomto\SheafHom_{\cal{E}/-}(\cal{H},\cal{F})\times
\SheafHom_{\cal{E}/-}(\cal{H},\cal{G})
$$
by taking the sets of objects of the two sides, where the functor
written is the one which has as components the functors $h\mto
\pr_1\circ h$ and $h\mto \pr_2\circ h$ from the first side to the two
factors of the second. The reader is left to verify that one
thus indeed obtains an isomorphism (the analogous fact being true,
more generally, whenever one has a projective limit
of categories --- and not only in the case of a fibered
product).

Let us moreover recall that one has (as was said in~no.~\ref{VI.1}):
\begin{eqnarray*}
\Ob(\cal{F}\times_\cal{E}\cal{G})&=&
\Ob(\cal{F})\times_{\Ob(\cal{E})}\Ob(\cal{G})\\
\Fl(\cal{F}\times_\cal{E}\cal{G})&=&
\Fl(\cal{F})\times_{\Fl(\cal{E})}\Fl(\cal{G})\quoi,
\end{eqnarray*}
the composition of arrows being performed moreover
componentwise.

In the sequel, we consider a functor
$$
\lambda\colon\cal{E}'\to\cal{E}
$$
and for every category $\cal{F}$ over~$\cal{E}$, one
considers $\cal{F}\times_\cal{E}\cal{E}'$ as a category
over~$\cal{E}'$ by means of $\pr_2$; in other words,
we interpret the operation ``fibered product'' as an
operation \emph{``change of base''}, the functor
$\lambda\colon\cal{E}'\to\cal{E}$ taking the name of \emph{``change of
base functor''.}
In accordance with the well-known general facts, one thus obtains
a functor, called the
% original p. 154
\emph{change of base functor} for $\lambda$:
$$
\lambda^*\colon\Cat_{/\cal{E}}\to\Cat_{/\cal{E}'}\quoi,
$$
(adjoint of the functor ``restriction of the base'' which, to every
category $\cal{F}'$ over~$\cal{E}'$, with
projection functor~$p'$, associates~$\cal{F}'$, regarded as a
category over~$\cal{E}$ by the functor~$p=\lambda
p'$). As is well known in the general case of a
change of base functor in a category, the
change of base functor ``commutes with projective limits'', and in
particular ``transforms'' fibered products over~$\cal{E}$ into
fibered products over~$\cal{E}'$.

Let $\cal{F}$ and $\cal{G}$ be two categories over
$\cal{E}$, one will define a \emph{canonical isomorphism}:
\begin{multline*}
\tag{i}
\label{eq:VI.3.i}
\SheafHom_{\cal{E}'/-}(\cal{F}',\cal{G}')\isomto
\SheafHom_{\cal{E}/-}(\cal{F}\times_{\cal{E}}\cal{E}',\cal{G})\\
\text{where }\cal{F}'=\cal{F}\times_{\cal{E}}\cal{E}'\text{, }
\cal{G}'=\cal{G}\times_{\cal{E}}\cal{E}'.
\end{multline*}
For this, consider the functor
$$
\pr_1\colon\cal{G}'=\cal{G}\times_{\cal{E}}\cal{E}'\to \cal{G}\quoi,
$$
and define~\eqref{eq:VI.3.i} by
$$
F\mto \pr_1\circ F\quoi,
$$
which a priori designates a functor
\begin{equation*}
\tag{ii}
\label{eq:VI.3.ii} {\SheafHom(\cal{F}',\cal{G}')\to\SheafHom(\cal{F}',\cal{G})}
\end{equation*}
It is therefore only necessary to verify that the latter induces a functor
for the subcategories~\eqref{eq:VI.3.i}, and that the latter is
an isomorphism. That~\eqref{eq:VI.3.ii} induces a bijection
$$
\Hom_{\cal{E}'/-}(\cal{F}',\cal{G}')\isomto
\Hom_{\cal{E}/-}(\cal{F}\times_{\cal{E}}\cal{E}',\cal{G})
$$
is the characteristic property of the change of
base functor. It remains therefore to
% original p. 155
prove that if $F$, $G$ are $\cal{E}'$-functors
$\cal{F}'\to\cal{G}'$, then \emph{the map}
$$
u\mto \pr_1\circ u
$$
\emph{induces a bijection}
$$
\Hom_{\cal{E}'}(F,G)\isomto\Hom_{\cal{E}}(\pr_1\circ F,\pr_1\circ
G)\quoi.
$$
The verification of this fact is immediate, and left to the
reader.

It follows from this isomorphism~\eqref{eq:VI.3.i}, and from the end of
the preceding no., that
$$
\SheafHom_{\cal{E}'/-}(\cal{F}\times_{\cal{E}}\cal{E}',\cal{G}
\times_{\cal{E}}\cal{E}')
$$
\emph{can be regarded as a functor in}
$\cal{E}',\cal{F},\cal{G}$, \emph{from the category}
$\Cat_{/\cal{E}}^\circ\times\Cat_{/\cal{E}}^\circ\allowbreak\times\Cat_{/\cal{E}}$
\emph{to the category}~$\Cat$, isomorphic to the functor defined
by the expression
$\SheafHom_{\cal{E}/-}(\cal{F}\times_{\cal{E}}\cal{E}',\cal{G})$. In
particular, for $\cal{F}$,~$\cal{G}$ fixed, one obtains a
functor in $\cal{E}'$, $\cal{E}'\to
\SheafHom_{\cal{E}''/-}(\cal{F}',\cal{G}')=
\SheafHom_{\cal{E}'/-}(\cal{F}\times_{\cal{E}}
\cal{E}',\cal{G}\times_{\cal{E}}\cal{E}')$, and in particular the
projection $\cal{E}$-functor $\lambda\colon \cal{E}'\to\cal{E}$
defines a morphism \ie a functor
$$
\lambda^*_{\cal{F},\cal{G}}\colon
\SheafHom_{\cal{E}/-}(\cal{F},\cal{G})\to
\SheafHom_{\cal{E}'/-}(\cal{F}',\cal{G}')
$$
which we shall make explicit. For the sets of objects of the two
sides, it is the map
$$
f\mto f'=f\times_{\cal{E}}\cal{E}'
$$
which expresses the functorial dependence of
$\cal{F}\times_{\cal{E}}\cal{E}'$ on the object $\cal{F}$
over~$\cal{E}$. On the other hand, consider two $\cal{E}$-functors
$$
f,g\colon\cal{F}\to \cal{G}
$$
and a homomorphism of~$\cal{E}$-functors
$$
u\colon f\to g\quoi,
$$
we
% original p. 156
shall make explicit the corresponding homomorphism of~$\cal{E}'$-functors:
$$
u'\colon f'\to g'\quoi.
$$
For every
$$
\xi'=(\xi,S')\in\Ob(\cal{F}')
$$
with
$$
\xi\in\Ob(\cal{F}),\quad S'\in\Ob(\cal{E}'),\quad p(\xi)=\lambda(S')=S
$$
the morphism
$$
u'(\xi')\colon f'(\xi')=(f(\xi),S')\to g'(\xi')=(g(\xi),S')\quad
\text{in } \cal{G}'
$$
is defined by the formula
$$
u'(\xi')=(u(\xi),\id_{S'})
$$
(which is indeed an $S'$-morphism in~$\cal{G}'$, because
$q(u(\xi))=\lambda(\id_{S'})=\id_S$).

Consider now an arbitrary $\cal{E}$-functor
$$
\lambda'\colon\cal{E}''\to\cal{E}'
$$
and the corresponding functor
$$
\SheafHom_{\cal{E}'/-}(\cal{F}\times_{\cal{E}}\cal{E}',
\cal{G}\times_{\cal{E}}\cal{E}')\to
\SheafHom_{\cal{E}''/-}(\cal{F}\times_{\cal{E}}\cal{E}'',
\cal{G}\times_{\cal{E}}\cal{E}'')\quoi,
$$
I say that this functor is nothing other than the functor one obtains by
the preceding procedure, starting from~$\cal{F}'$ and
$\cal{G}'$ over~$\cal{E}'$ and regarding $\cal{E}''$ as a
category over~$\cal{E}'$, taking into account the isomorphisms of
\emph{``transitivity of change of base''}:
$$
\cal{F}'\times_{\cal{E}'}\cal{E}''\isomto
\cal{F}''=\cal{F}\times_{\cal{E}}\cal{E}''\quad
\text{and}\quad\cal{G}'\times_{\cal{E}'}\cal{E}''\isomto\cal{G}''
=\cal G\times_{\cal E}\cal E''\quoi,
$$
which imply a canonical isomorphism
$$
\SheafHom_{\cal{E}''/-}(\cal{F}'\times_{\cal{E}'}\cal{E}'',
\cal{G}'\times_{\cal{E}'}\cal{E}'')\isomto
\SheafHom_{\cal{E}''/-}(\cal{F}\times_{\cal{E}}\cal{E}'',
\cal{G}\times_{\cal{E}}\cal{E}'')
$$
The
% original p. 157
verification of this compatibility is immediate, and
left to the reader.

The functors just defined are compatible with the
pairings defined in the preceding no., in a precise way: if $\cal{F}$, $\cal{G}$,~$\cal{H}$ are
categories over~$\cal{E}$ and if one sets
$$
\cal{F}'=\cal{F}\times_{\cal{E}}\cal{E}'\quoi,\quad \cal{G}'=
\cal{G}\times_{\cal{E}}\cal{E}'\quoi,\quad \cal{H}'=
\cal{H}\times_{\cal{E}}\cal{E}'\quoi,
$$
one has commutativity in the following diagram of functors:
$$
\xymatrix{ \SheafHom_{\cal{E}/-}(\cal{F},\cal{G})\times
\SheafHom_{\cal{E}/-}(\cal{G},\cal{H}) \ar[r]
\ar[d]_{\lambda^*_{\cal{F},\cal{G}}\times \lambda^*_{\cal{G},\cal{H}}}
& \SheafHom_{\cal{E}/-}(\cal{F},\cal{H})
\ar[d]^{\lambda^*_{\cal{F},\cal{H}}} \\
\SheafHom_{\cal{E}'/-}(\cal{F}',\cal{G}')\times
\SheafHom_{\cal{E}'/-}(\cal{G}',\cal{H}')\ar[r] &
\SheafHom_{\cal{E}'/-}(\cal{F}',\cal{H}')}
$$
where the horizontal arrows are the composition functors
defined in the preceding no. This commutativity
is expressed by the formulae
$$
(gf)'=g'f'
$$
for $f\in\Hom_{\cal{E}}(\cal{F},\cal{G})$,
$g\in\Hom_{\cal{E}}(\cal{G},\cal{H})$, (a formula which simply expresses
the functoriality of change of base), and the formula
$$
(v*u)'=v'*u'
$$
when $u\colon f\to f_1$ is an arrow of
$\SheafHom_{\cal{E}/-}(\cal{F},\cal{G})$ and $v\colon g\to g_1$ an
arrow of~$\SheafHom_{\cal{E}/-}(\cal{G},\cal{H})$. The
verification of this formula follows easily from the
definitions.

In the sequel, we shall be interested mainly in
$\SheafHom_{\cal{E}}(\cal{F},\cal{G})$ (and certain
remarkable subcategories of the latter) when
$\cal{F}=\cal{E}$, and we introduce for this reason a
special notation:
$$
\mathbf{\Gamma}(\cal{G}/\cal{E})=
\SheafHom_{\cal{E}}(\cal{E},\cal{G})\quoi,\quad
\Gamma(\cal{G}/\cal{E})=\Ob(\mathbf{\Gamma}(\cal{G}/\cal{E}))=
\Hom_{\cal{E}}(\cal{E},\cal{G})\quoi.
$$

\begin{remarks}
When
% original p. 158
$\cal{E}$ is a punctual category,
\ie $\Ob(\cal{E})$ and $\Fl(\cal{E})$ reduced to a single
element, which also means that $\cal{E}$ is a
final object of the category~$\Cat$, then the datum of a
category over~$\cal{E}$ is equivalent to the datum
of just a category, (because there will be a unique functor from
$\cal{F}$ to~$\cal{E}$). More precisely,
$\Cat_{/\cal{E}}$ is then isomorphic to~$\Cat$. Moreover, the
categories $\SheafHom_{\cal{E}/-}(\cal{F},\cal{G})$ are then nothing
other than the $\SheafHom(\cal{F},\cal{G})$. Let us then recall that the
fundamental formula
$$
\Hom(\cal{H},\SheafHom(\cal{F},\cal{G}))\isomto
\Hom(\cal{F}\times\cal{H},\cal{G})
$$
(a functorial isomorphism in the three arguments which appear there),
allows one to interpret $\SheafHom(\cal{F},\cal{G})$ axiomatically,
in terms internal to the category~$\Cat$, so that the
known formulary of the $\SheafHom$ categories appears as a
special case of a formulary valid in categories such
as~$\Cat$, where ``$\SheafHom$-objects'' (defined by the
preceding formula) exist. There is an analogous
interpretation of~$\SheafHom_{\cal{E}/-}(\cal{F},\cal{G})$ when one assumes
once more that $\cal{E}$ is arbitrary, by the formula
$$
\Hom(\cal{H},\SheafHom_{\cal{E}/-}(\cal{F},\cal{G}))\isomto
\Hom_{\cal{E}}(\cal{F}\times\cal{H},\cal{G})
$$
(a functorial isomorphism in the three arguments). In this way,
the formal properties set out in~no.~\ref{VI.2},~\ref{VI.3} are
special cases of more general results, valid
in categories where the objects
$\SheafHom_{\cal{E}/-}(\cal{F},\cal{G})$ (when $\cal{F}$,~$\cal{G}$
are two objects of the category over a third~$\cal{E}$) exist.
\end{remarks}

% ---- en-chunk2.tex ----

\section[Fiber-categories]{Fiber-categories; equivalence of
$\cal{E}$-categories}
\label{VI.4}

Let~$\cal{F}$ be a category over~$\cal{E}$, and let $S \in
\Ob({\cal{E}})$. One calls \emph{fiber-category of}~$\cal{F}$
at~$S$
the subcategory~$\cal{F}_S$
of~$\cal{F}$, inverse image of the
% original p. 159
punctual subcategory of~$\cal{E}$ defined by $S$. Thus
the objects of~$\cal{F}_S$ are the objects~$\xi$ of~$\cal{F}$ such that
$p(\xi)=S$, its morphisms are the morphisms~$u$ of~$\cal{F}$ such
that $p(u)=\id_S$, \ie the $S$-morphisms in $\cal{F}$. Of
course,~$\cal{F}_S$ is canonically isomorphic to the fibered product
$\cal{F} \times_{\cal{E}} \{S\}$, where~$\{S\}$ denotes the
punctual subcategory of~$\cal{E}$ defined by~$S$, equipped
with its inclusion functor into~$\cal{E}$. It follows (taking
into account the transitivity of change of base) that if one performs the
change of base $\lambda\colon \cal{E}' \to \cal{E}$, then for
every $S' \in \Ob(\cal{E}')$, \emph{the projection $\pr_1\colon
\cal{F}' = \cal{F} \times_{\cal{E}} \cal{E}' \to \cal{F}$ induces an
isomorphism}
$$
\cal{F}'_{S'} \to \cal{F}_S \qquad (\text{where $S=\lambda(S')$}).
$$

\begin{proposition}
\label{VI.4.1} Let $f\colon \cal{F} \to \cal{G}$ be an
$\cal{E}$-functor. If~$f$ is fully faithful, then for every
change of base $\cal{E}' \to \cal{E}$, the corresponding functor
$f' \colon \cal{F}' = \cal{F} \times_{\cal{E}} \cal{E}' \to\allowbreak \cal{G}' =
\cal{G} \times_{\cal{E}} \cal{E}'$ is fully faithful.
\end{proposition}

The verification is immediate; more generally, one
can show that every projective limit of
fully faithful functors
(here,~$f$ and the identity functors
in~$\cal{E},\cal{E}'$) is a fully faithful functor.

One notes that the assertion analogous to 4.1,
where
``fully faithful'' is replaced by ``equivalence
of categories'', is false, already for~$\cal{G}=\cal{E}$.
However:

\begin{proposition}
\label{VI.4.2} Let $f \colon \cal{F} \to \cal{G}$ be an
$\cal{E}$-functor. The following conditions are equivalent:
\begin{enumerate}
\item[(i)] There exists an $\cal{E}$-functor $g \colon \cal{G} \to
\cal{F}$ and $\cal{E}$-isomorphisms
$$
gf \isomto \id_{\cal{F}}, \quad fg \isomto \id_{\cal{G}}.
$$
\item[(ii)] For every category $\cal{E}'$ over~$\cal{E}$, the
functor
$$
f' = f \times_{\cal{E}} \cal{E}' \colon \cal{F}' = \cal{F}
\times_{\cal{E}} \cal{E}' \to \cal{G}'=\cal{G} \times_{\cal{E}}
\cal{E}'
$$
is an equivalence of categories.
\item[(iii)]
% original p. 160
$f$ is an equivalence of categories, and for every $S \in
\Ob(\cal{E})$, the functor $f_S \colon \cal{F}_S \to \cal{G}_S$ induced
by~$f$ is an equivalence of categories.
\item[(iii~bis)] $f$ is fully faithful, and for every $S \in
\Ob(\cal{E})$ and every
$\eta\in\Ob(\cal{G}_S)$,
there exists a $\xi \in \Ob(\cal{F}_S)$ and an $S$-isomorphism
$u \colon f(\xi) \to \eta$.
\end{enumerate}
\end{proposition}

\subsubsection*{Proof} Obviously (i) implies that~$f$ is an
equivalence of categories (a notion which is defined by the
same condition, but where the isomorphisms of functors are
not required to be $\cal{E}$-morphisms). On the other hand,
it follows from the functorialities of the preceding no.
that condition (i) is preserved after change of
base~$\cal{E}'\to\cal{E}$. It follows that (i)$\To$(ii).
Obviously (ii)$\To$(iii), for it suffices to take $\cal{E}'
= \cal{E}$ and~$\cal{E}'=\{S\}$. It is even more trivial that
(iii)$\To$(iii~bis), it remains to prove that (iii~bis)$\To$(i). For this, let us choose for every $\eta \in
\Ob(\cal{G})$ a $g(\eta) \in \Ob(\cal{F})$ and an isomorphism
$u(\eta)\colon f(g(\eta))\to\eta$ which is such that~$q(u(\eta)) =
\id_S$, where~$S=q(\eta)$. This is possible thanks to the
second condition (iii~bis). As is well-known and immediate,
the fact that~$f$ is fully faithful implies that~$g$ can in a
unique way be regarded as a functor
from~$\cal{G}$ to~$\cal{F}$, in such a way that the~$u(\eta)$
define a functorial homomorphism (hence an isomorphism)
$u\colon fg \isomto \id_{\cal{G}}$. Moreover, by construction~$g$ is
an $\cal{E}$-functor and~$u$ an $\cal{E}$-homomorphism. To the
preceding data there then corresponds a functorial isomorphism
$v\colon gf \to \id_{\cal{F}}$, defined by the condition
that $f*v=u*f$, and one sees at once that it is likewise an
$\cal{E}$-morphism, qed.

\begin{definition}
\label{VI.4.3} If the preceding conditions are
satisfied, one says that~$f$ is an equivalence of
categories over~$\cal{E}$, or an $\cal{E}$-equivalence.
\end{definition}

\begin{corollary}
\label{VI.4.4} Suppose that the projection functor $p\colon
\cal{F}\to\cal{E}$ is a transportable functor, \ie that for every
isomorphism $\alpha\colon T \to S$ in~$\cal{E}$ and every object~$\xi$
in~$\cal{F}_T$, there exists an isomorphism~$u$ in~$\cal{F}$ with
source~$\xi$ such that~$p(u)=\alpha$. Then every $\cal{E}$-functor
$f\colon \cal{F}\to\cal{G}$ which is an equivalence of
categories, is an $\cal{E}$-equivalence.
\end{corollary}

Follows
% original p. 161
from the criterion (iii~bis).

\begin{corollary}
\label{VI.4.5} Let $f\colon \cal{F}\to\cal{G}$ be an
$\cal{E}$-equivalence. Then for every category~$\cal{H}$
over~$\cal{E}$, the corresponding functors:
\begin{align*}
\SheafHom_{\cal{E}/-}(\cal{G},\cal{H}) & \to
\SheafHom_{\cal{E}/-}(\cal{F},\cal{H}) \\
\SheafHom_{\cal{E}/-}(\cal{H},\cal{F})
& \to \SheafHom_{\cal{E}/-}(\cal{H},\cal{G})
\end{align*}
(\cf no.~\ref{VI.2}) are equivalences of categories.
\end{corollary}

This follows from the criterion (i) by the usual reasoning.

\section{Cartesian morphisms, inverse images, cartesian functors}
\label{VI.5}

Let~$\cal{F}$ be a category over~$\cal{E}$, with
projection functor~$p$.

\begin{definition}
\label{VI.5.1} Consider a morphism
$$
\alpha\colon \eta\to\xi
$$
in~$\cal{F}$, and let
$$
S = p(\xi), \quad T = p(\eta), \quad f=p(\alpha).
$$
One says that~$\alpha$ is a \emph{cartesian morphism}
if for every $\eta' \in \Ob(\cal{F}_T)$ and every $f$-morphism $u\colon
\eta'\to\xi$, there exists a unique $T$-morphism $\overline{u}\colon
\eta'\to\eta$ such that~$u=\alpha\circ\overline{u}$.
\end{definition}

This therefore means that for every $\eta'\in\Ob(\cal{F}_T)$,
the map $v \mto \alpha\circ v$:
\begin{equation}
\tag{i} \Hom_T(\eta',\eta) \to \Hom_f(\eta',\xi)
\end{equation}
is bijective. This also means that the pair $(\eta,\alpha)$
\emph{represents the functor in~$\eta'$} $\cal{F}^{\circ}_T \to
\Ens$ of the right-hand side. If for a given morphism $f\colon T \to S$
in~$\cal{E}$, and a given $\xi \in \Ob(\cal{F}_S)$, there
exists such a pair $(\eta,\alpha)$, \ie a
% original p. 162
cartesian morphism~$\alpha$ in~$\cal{F}$ with target~$\xi$, such that
$p(\alpha)=f$, then~$\eta$ is determined in $\cal{F}_T$
up to unique isomorphism. One then says that \emph{the inverse
image of~$\xi$ by~$f$ exists}, and an object~$\eta$ of~$\cal{F}_T$
equipped with a cartesian $f$-morphism $\alpha\colon \eta \to \xi$ is
called \emph{an inverse image of~$\xi$ by~$f$}.
Often one assumes
such an inverse image chosen whenever it exists ($\cal{F}$
being fixed); one then denotes the inverse image by symbols
such as~$f^*_{\cal{F}}(\xi)$, or simply~$f^*(\xi)$
or~$\xi \times_S T$
when these notations do not lead to
confusion;
the canonical morphism $\alpha\colon \eta \to \xi$ will then be
denoted, in what follows, by~$\alpha_f(\xi)$.
If for every $\xi \in \Ob(\cal{F}_S)$, the inverse image of~$\xi$
by~$f$ exists, one will also say that \emph{the inverse image functor
by~$f$ in~$\cal{F}$ exists}, and~$f^*(\xi)$ then becomes a
\emph{functor covariant in~$\xi$}, from~$\cal{F}_S$ to~$\cal{F}_T$.
This comes from the fact that the right-hand side in (i) depends
covariantly on~$\xi$, \ie more precisely
designates a functor from~$\cal{F}^{\circ}_T \times \cal{F}_S$
to~$\Ens$. This functorial dependence for~$f^*(\xi)$
is made explicit as follows: consider cartesian $f$-morphisms
$$
\alpha\colon \eta \to \xi, \quad \alpha'\colon \eta' \to \xi'
$$
and an $S$-morphism $\lambda\colon \xi \to \xi'$, then there exists a
unique $T$-morphism $\mu\colon \eta\to\eta'$ such that one has
$$
\alpha' \mu = \lambda \alpha
$$
(as follows from the fact that~$\alpha'$ is cartesian).

Let us also note the following immediate fact: consider a
commutative diagram
$$
\xymatrix@C=1.2cm{
\xi\ar[d]_-{\lambda}& \ar[l]_-{\alpha}\eta \ar[d]^-{\mu}\\
\xi' &\ar[l]_-{\alpha'}\eta'
}
$$
in~$\cal{F}$,
% original p. 163
where~$\alpha$ and~$\alpha'$ are $f$-morphisms, and~$\lambda$ an
$S$-isomorphism,~$\mu$ a $T$-isomorphism. \emph{For~$\alpha$
to be cartesian, it is necessary and sufficient that~$\alpha'$
be so}.

\begin{definition}
\label{VI.5.2} An $\cal{E}$-functor $F\colon \cal{F}\to\cal{G}$ is
called a \emph{cartesian functor}
if it transforms cartesian morphisms into cartesian
morphisms. One denotes by $\SheafHom_\cart(\cal{F},\cal{G})$
the full subcategory of
$\SheafHom_{\cal{E}/-}(\cal{F},\cal{G})$ formed of the
cartesian functors.

For example, considering~$\cal{E}$ as a category
over~$\cal{E}$ by means of the identity functor, every morphism
of~$\cal{E}$ is cartesian, hence a cartesian functor
from~$\cal{E}$ to~$\cal{F}$ is a section functor $F\colon
\cal{E}\to\cal{F}$ which transforms every morphism of~$\cal{E}$ into a
cartesian morphism; such a functor is called a \emph{cartesian
section} of~$\cal{F}$ over~$\cal{E}$.
\end{definition}

\begin{proposition}
\label{VI.5.3} \textup{(i)}~A functor $F\colon \cal{F} \to \cal{G}$ which is an
$\cal{E}$-equivalence, is a cartesian functor.
\textup{(ii)}~Let~$F,G$ be two \emph{isomorphic} $\cal{E}$-functors~$\cal{F}\to\cal{G}$. If one is cartesian,
so is the other. \textup{(iii)}~The composite of two cartesian functors
$\cal{F}\to\cal{G}$ and $\cal{G}\to\cal{H}$ is a
cartesian functor.
\end{proposition}

Assertion (iii) is trivial from the definition, (ii)
follows from the remark preceding~\ref{VI.5.2}, (i) follows
easily from the definition and from the criterion~\ref{VI.4.2}~(iii); more
precisely, a morphism~$\alpha$ in~$\cal{F}$ is
cartesian if and only if~$F(\alpha)$ is so.

\begin{corollary}
\label{VI.5.4}
Let $F \colon \cal{F} \to \cal{G}$ be an
$\cal{E}$-equivalence. Then for every category~$\cal{H}$
over~$\cal{E}$, the corresponding functors $G \mto G \circ F$ and
$G \mto F \circ G$ induce equivalences of
categories:
\begin{eqnarray*}
\SheafHom_\cart(\cal{G},\cal{H})\lto{\approx}
\SheafHom_\cart(\cal{F},\cal{H}) \\
\SheafHom_\cart(\cal{H},\cal{F})\lto{\approx}
\SheafHom_\cart(\cal{H}, \cal{G})
\end{eqnarray*}
\end{corollary}

This is deduced in the usual way from~\ref{VI.4.2}
criterion~(i) and from~\ref{VI.5.3} (i) (ii)~(iii).
% original p. 164
One can make more precise that \emph{the $\cal{E}$-functor
$G\colon \cal{G} \to \cal{H}$ is cartesian if and only if $G
\circ F$ is so}, and likewise \emph{an $\cal{E}$-functor $G\colon
\cal{H} \to \cal{F}$ is cartesian if and only if $F \circ G$
is so}.

It follows from~\ref{VI.5.4}~(iii) that if one considers the
subcategory~$\Cat^\cart_{/\cal{E}}$
of~$\Cat_{/\cal{E}}$ whose objects are the same as those of
$\Cat_{/\cal{E}}$, and whose morphisms are the
\emph{cartesian} functors, then one has as in~no.~\ref{VI.2} pairings:
$$
\SheafHom_\cart(\cal{F},\cal{G})\times
\SheafHom_\cart(\cal{G},\cal{H}) \to \SheafHom_\cart(\cal{F},\cal{H})
$$
induced
by those of~no.~\ref{VI.2}, allowing one to consider
$\SheafHom_\cart(\cal{F},\cal{G})$ as a functor in $\cal{F}$,
$\cal{G}$, from the category $\left(\Cat^\cart_{/\cal{E}}
\right)^\circ \times \Cat^\cart_{/\cal{E}}$
to~$\Cat$.
We will need this remark especially for the case where
$\cal{F} = \cal{G}$:

\begin{definition}
\label{VI.5.5}
Let $\cal{F}$ be a category over~$\cal{E}$. One denotes by
$$
\varprojLim\cal{F}/\cal{E}
$$
the category of cartesian $\cal{E}$-functors
$\cal{E}\to\cal{F}$, \ie of cartesian sections of~$\cal{F}$
over~$\cal{E}$.
\end{definition}

According to what has just been said,
$\varprojLim\cal{F}/\cal{E}$ is a functor
in~$\cal{F}$, from the category $\Cat^\cart_{/\cal{E}}$ to the
category $\Cat$.

We will see below the relations between this operation
$\varprojLim$ and the notion of projective limit of
categories, as well as numerous examples.

\section[Fibered categories and prefibered categories]{Fibered categories and prefibered categories. Products and change of base therein}
\label{VI.6}

\begin{definition}
\label{VI.6.1}
A category $\cal{F}$ over~$\cal{E}$ is called a
\emph{fibered category}
(and one then says that the functor $\cal{F}\to\cal{E}$ is
\emph{fibrant})
if it satisfies the following two axioms:
\begin{itemize}
\item{$\Fib_{\mathrm{I}}$}
% original p. 165
For every morphism $f\colon T\to S$ in~$\cal{E}$, the inverse image functor
by $f$ in $\cal{F}$ exists.
\item{$\Fib_{\mathrm{II}}$} The composite of two cartesian morphisms
is cartesian.
\end{itemize}
A category $\cal{F}$ over~$\cal{E}$ satisfying the condition
$\Fib_{\mathrm{I}}$ is called a \emph{prefibered category
over~$\cal{E}$}.
\end{definition}

If~$\cal{F}$ is a fibered category
(\resp prefibered) over~$\cal{E}$, a subcategory
$\cal{G}$ of~$\cal{F}$ is called a \emph{fibered subcategory}
(\resp a \emph{prefibered subcategory}) if it is a
fibered category (\resp prefibered) over~$\cal{E}$, and
if moreover the inclusion functor is cartesian. If for example
$\cal{G}$ is a \emph{full} subcategory of~$\cal{F}$, one
sees that this means that for every morphism $f\colon T\to S$ in
$\cal{E}$ and for every $\xi \in \Ob(\cal{G}_S)$, $f^*_{\cal{F}}(\xi)$
is $T$-isomorphic to an object of~$\cal{G}_T$. Another
interesting case is the following: $\cal{F}$ being a
fibered category over~$\cal{E}$, consider the subcategory
$\cal{G}$ of~$\cal{F}$ having the same objects, and whose morphisms
are the \emph{cartesian} morphisms of~$\cal{F}$; in particular
the morphisms of~$\cal{G}_S$ are the isomorphisms of
$\cal{F}_S$. One sees at once that this is indeed a fibered subcategory
of~$\cal{F}$, for in the bijection
$$
\Hom_T(\eta',\eta) \isomto \Hom_f(\eta',\xi)
$$
relative to a cartesian $f$-morphism $\alpha$ in $\cal{F}$,
to the $T$-isomorphisms of the left-hand side correspond the cartesian
morphisms of the right-hand side. By definition, the cartesian
sections $\cal{E}\to\cal{F}$
then correspond bijectively to arbitrary $\cal{E}$-functors
$\cal{E}\to\cal{G}$ (but one notes that the natural functor
$$
\SheafHom_{\cal{E}/-}(\cal{E},\cal{G})\to
\SheafHom_\cart(\cal{E},\cal{F})=
\varprojLim(\cal{F}/\cal{E})
$$
is faithful, but in general is not fully
faithful, \ie is not an isomorphism).

\begin{remarks}
Let $\cal{F}$ be a category over~$\cal{E}$. The following
conditions are equivalent:~(i)~All morphisms of~$\cal{F}$
are cartesian (ii)~$\cal{F}$~is a fibered category
over~$\cal{E}$, and the $\cal{F}_S$ are groupoids
(\ie every morphism in $\cal{F}_S$ is an isomorphism). One
then says that $\cal{F}$ is a category \emph{fibered in
groupoids}
% original p. 166
over~$\cal{E}$. These are the ones one encounters especially in ``theory
of modules''. If $\cal{E}$ is a groupoid, one shows that
conditions (i) and (ii) are also equivalent to the following:
(iii)~$\cal{F}$ is a groupoid, and the projection functor
$p\colon \cal{F}\to\cal{E}$ is transportable (\cf \ref{VI.4.4}). For
example, if $\cal{E}$ and $\cal{F}$ are groupoids such that
$\Ob(\cal{E})$ and $\Ob(\cal{F})$
are reduced to a point, so that $\cal{E}$ and $\cal{F}$
are defined, up to isomorphism, by groups $E$ and
$F$, and the functor $p\colon \cal{F}\to\cal{E}$ is defined by a
group homomorphism $p\colon F\to E$, then $\cal{F}$ is
fibered over~$\cal{E}$ if and only if $p$ is surjective, \ie if
$p$ defines an \emph{extension} of the group $E$ by the group $G =
\Ker p$.
\end{remarks}

\begin{proposition}
\label{VI.6.2}
Let $F \colon \cal{F} \to \cal{G}$ be an
$\cal{E}$-equivalence. For $\cal{F}$ to be a
fibered category (\resp prefibered) over~$\cal{E}$, it is necessary and
sufficient that $\cal{G}$ be so.
\end{proposition}

Follows easily from the definitions and from the remark
noted above that a morphism
$\alpha$
in $\cal{F}$ is cartesian if and only if $F(\alpha)$ is so.

\begin{proposition}
\label{VI.6.3}
Let $\cal{F}_1$, $\cal{F}_2$ be two categories over~$\cal{E}$, and
let $\alpha = (\alpha_1,\alpha_2)$ be a morphism in
$\cal{F}=\cal{F}_1 \times_{\cal{E}} \cal{F}_2$. For $\alpha$ to be
cartesian, it is necessary and sufficient that its components be so.
\end{proposition}

Indeed, let $\xi_i$ be the target and $\eta_i$ the source of~$\alpha_i$, and
let $f\colon T\to S$ be the morphism of~$\cal{E}$ such that $\alpha_1$ and
$\alpha_2$ are $f$-morphisms. For every
$\eta'=(\eta'_1,\eta'_2)$ in $\cal{F}_T$, one has a commutative
diagram
$$
\xymatrix{ \Hom_T(\eta',\eta) \ar[r] \ar[d] & \Hom_f(\eta',\xi) \ar[d]
\\ \Hom_T(\eta'_1,\eta_1)\times\Hom_T(\eta'_2,\eta_2) \ar[r] &
\Hom_f(\eta'_1,\xi_1)\times \Hom_f(\eta'_2,\xi_2) }
$$
where the vertical arrows are bijections. Thus if one
of the horizontal arrows is a bijection, so is
the other. This already shows that if $\alpha_1$, $\alpha_2$
are cartesian (hence the second horizontal arrow
bijective) then $\alpha$ is so. The converse is seen by setting
in the diagram above $\eta'_i=\eta_i$ whence
$\Hom_T(\eta'_i,\eta_i) \neq \emptyset$, first for $i=2$ which
proves that $\alpha_1$ is cartesian, then for $i=1$ which proves
that $\alpha_2$ is cartesian.

\begin{corollary}
\label{VI.6.4}
Let
% original p. 167
$\cal{F} = \cal{F}_1\times_{\cal{E}}\cal{F}_2$, and let
$F=(F_1,F_2)$ be an $\cal{E}$-functor $\cal{G}\to\cal{F}$. For
$F$
to be cartesian, it is necessary and sufficient that $F_1$ and $F_2$
be so. One thus obtains an isomorphism of categories
$$
\SheafHom_\cart(\cal{G}, \cal{F}_1 \times_{\cal{E}} \cal{F}_2) \isomto
\SheafHom_\cart(\cal{G},\cal{F}_1)
\times\SheafHom_\cart(\cal{G},\cal{F}_2)
$$
and in particular (setting $\cal{G}=\cal{E}$) an isomorphism of
categories
$$
\varprojLim(\cal{F}_1 \times_{\cal{E}}
\cal{F}_2/\cal{E}) \isomto
\varprojLim(\cal{F}_1/\cal{E}) \times
\varprojLim(\cal{F}_2/\cal{E})
$$
\end{corollary}

\begin{corollary}
\label{VI.6.5}
Let $\cal{F}_1$ and $\cal{F}_2$ be two fibered categories
(\resp prefibered) over~$\cal{E}$, then their fibered
product $\cal{F}=\cal{F}_1 \times_{\cal{E}}\cal{F}_2$ is a
fibered category (\resp prefibered) over~$\cal{E}$.
\end{corollary}

These results moreover extend to the case of the fibered
product of an arbitrary family of categories over~$\cal{E}$.

\begin{proposition}
\label{VI.6.6}
Let $\cal{F}$ be a category over~$\cal{E}$, with
projection functor $p$, and let $\lambda \colon \cal{E}' \to \cal{E}$
be a functor, consider
$\cal{F}'=\cal{F}\times_{\cal{E}}\cal{E}'$ as a category over~$\cal{E}'$ by the projection functor $p'=p \times_{\cal{E}}
\id_{\cal{E}'}$. Let $\alpha'$ be a morphism of~$\cal{F}'$; for
$\alpha'$ to be a cartesian morphism, it is necessary and sufficient that
its image $\alpha$ in $\cal{F}$ be so.
\end{proposition}

The proof is immediate and left to the reader.

\begin{corollary}
\label{VI.6.7}
For every cartesian functor $F\colon \cal{F} \to \cal{G}$ of
categories over~$\cal{E}$, the functor
$F'=F\times_{\cal{E}}\cal{E}'$ from
$\cal{F}'=\cal{F}\times_{\cal{E}}\cal{E}'$ to
$\cal{G}'=\cal{G}\times_{\cal{E}}\cal{E}'$ is cartesian.
\end{corollary}

Consequently, the functor $\SheafHom_{\cal{E}}(\cal{F},\cal{G}) \to
\SheafHom_{\cal{E}'}(\cal{F}',\cal{G}')$ considered in
no.~\ref{VI.3} induces a functor
$$
\SheafHom_\cart(\cal{F},\cal{G}) \to
\SheafHom_{\cart}(\cal{F}',\cal{G}')\,;
$$
in other terms, for $\cal{F}$, $\cal{G}$ fixed, \emph{one can
consider}
$$
\SheafHom_\cart(\cal{F} \times_{\cal{E}} \cal{E}',
\cal{G}\times_{\cal{E}} \cal{E}')
$$
\emph{as
% original p. 168
a functor in~$\cal{E}'$, from the category
${\Cat_{/\cal{E}}}^\circ$ to $\Cat$}. If one likewise lets
$\cal{F}$, $\cal{G}$ vary, one finds
a functor from the category ${\Cat_{/\cal{E}}}^\circ \times \big(
\Cat^\cart_{/\cal{E}} \big)^\circ \times \Cat^\cart_{/\cal{E}}$
to~$\Cat$. When one takes into account the isomorphism
$$
\Hom_{\cal{E}'}(\cal{F}',\cal{G}') \isomto
\Hom_\cal{E}(\cal{F}\times_\cal{E}\cal{E}',\cal{G})
$$
envisaged in no.~\ref{VI.3}, then the cartesian $\cal{E}'$-functors
from~$\cal{F}'$ to~$\cal{G}'$ correspond to the $\cal{E}$-functors
$\cal{F} \times_\cal{E} \cal{E}' \to \cal{G}$ which transform every
morphism whose first projection is a cartesian
morphism of~$\cal{F}$, into a cartesian morphism of
$\cal{G}$. Setting $\cal{F}=\cal{E}$, one finds (after change
of notation):

\begin{corollary}
\label{VI.6.8}
$\varprojLim(\cal{F}'/\cal{E}')$ is isomorphic to
the full subcategory of
$\SheafHom_{\cal{E}/-}(\cal{E}',\cal{F})$ formed of the
$\cal{E}$-functors $\cal{E}'\to\cal{F}$ which transform arbitrary
morphisms into cartesian morphisms. In particular, if $\cal{F}$
is a fibered category and if $\tilde{\cal{F}}$ is the
subcategory of~$\cal{F}$ whose morphisms are the
cartesian morphisms of~$\cal{F}$, then one has a bijection
$$
\Ob \varprojLim(\cal{F}'/\cal{E}') \isomto
\Hom_{\cal{E}/-}(\cal{E}',\tilde{\cal{F}})\,.
$$
\end{corollary}

This makes more precise the way in which the expression
$\varprojLim(\cal{F} \times_{\cal{E}} \cal{E}' /
\cal{E}')$ is to be considered as a functor in
$\cal{E}'$ and in $\cal{F}$, from the category
${\Cat_{/\cal{E}}}^\circ \times \Cat^\cart_{/\cal{E}}$ to the
category $\Cat$. One will see later a more complete
functorial dependence with respect to $\cal{E}'$, when
$\cal{F}$ is required to be a fibered category.

\begin{corollary}
\label{VI.6.9}
Let $\cal{F}$ be a fibered category (\resp prefibered)
over~$\cal{E}$, then $\cal{F}'=\cal{F}\times_{\cal{E}}\cal{E}'$ is
a fibered category (\resp prefibered) over~$\cal{E}'$.
\end{corollary}

\begin{proposition}
\label{VI.6.10} Let $\cal{F}$, $\cal{G}$ be
prefibered categories over~$\cal{E}$, $F$ a cartesian $\cal{E}$-functor
from~$\cal{F}$ to $\cal{G}$. For $F$ to be
faithful,
% original p. 169
(\resp fully faithful, \resp an $\cal{E}$-equivalence)
it is necessary and sufficient that for every
$S\in \Ob(\cal{E})$,
the induced functor $F_S\colon \cal{F}_S\to \cal{G}_S$ be
faithful (\resp fully faithful, \resp an equivalence).
\end{proposition}

Immediate proof from the definitions.

To finish this no., we give some properties of
fibered categories, using the axiom $\Fib_{\mathrm{II}}$.
\begin{proposition}
\label{VI.6.11} Let $\cal{F}$ be a prefibered category
over~$\cal{E}$. For $\cal{F}$ to be fibered, it is necessary and
sufficient that it satisfy the following condition:

$\Fib_{\mathrm{II}}'$: Let $\alpha\colon\eta\to\xi$ be a
cartesian morphism in $\cal{F}$ over the morphism $f\colon T\to S$
of~$\cal{E}$. For every morphism $g\colon U\to T$ in $\cal{E}$, and
every
$\zeta\in\Ob(\cal{F}_U)$,
the map $u\mto\alpha\circ u$:
$$
\Hom_g(\zeta,\eta)\to\Hom_{fg}(\zeta,\xi)
$$
is bijective.
\end{proposition}

In other terms, in a \emph{fibered} category
over~$\cal{E}$, cartesian diagrams are characterized
by a property, a priori stronger than that of the
definition (which one obtains by setting $g=\id_T$ in
the preceding statement).
\begin{corollary}
\label{VI.6.12} Let $\cal{F}$ be a category over~$\cal{E}$,
$\alpha$ a morphism in $\cal{F}$. For $\alpha$ to be an
isomorphism, it is necessary that $p(\alpha)=f$ be an isomorphism and that
$\alpha$ be cartesian; the converse is true if $\cal{F}$
is fibered over~$\cal{E}$.
\end{corollary}

Indeed, if $\alpha$ is an isomorphism then so is evidently
$f=p(\alpha)$; for every
$\eta'\in\Ob(\cal{F}_T)$,
the map $u\mto \alpha\circ u$
$$
\Hom(\eta',\eta)\to \Hom(\eta',\xi)
$$
is bijective; as $f$ is an isomorphism, one sees at once
that an
element of the left-hand side is a $T$-morphism if and only
if its image in the right-hand side is an $f$-morphism, hence one thus obtains a bijection
$$
\Hom_T(\eta',\eta)\to\Hom_f(\eta',\xi)
$$
which
% original p. 170
proves the first assertion. Conversely, suppose
that $f$ is an isomorphism and that $\alpha$ satisfies the condition
stated in $\Fib_{\mathrm{II}'}$ (which therefore means,
when $\cal{F}$ is
fibered
over~$\cal{E}$, that $\alpha$ is cartesian), then one sees at
once that for every $\zeta\in\Ob\cal{F}$, the map $u\mto
\alpha\circ u$ from~$\Hom(\zeta,\eta)$ to $\Hom(\zeta,\xi)$ is
bijective, hence $\alpha$ is an isomorphism.
\begin{corollary}
\label{VI.6.13} Let $\alpha\colon\eta\to\xi$ and
$\beta\colon\zeta\to \eta$ be two composable morphisms in the
category $\cal{F}$ fibered over~$\cal{E}$. If $\alpha$ is
cartesian then $\beta$ is so if and only if $\alpha\beta$
is so.
\end{corollary}

One uses the definition of cartesian morphisms in the
strengthened form of~\ref{VI.6.11}.

% ---- en-chunk3.tex ----

\section{Cloven categories over~$\cal{E}$}
\label{VI.7}

\begin{definition}
\label{VI.7.1} Let $\cal{F}$ be a category over~$\cal{E}$. One
calls \emph{cleavage}
of~$\cal{F}$ over~$\cal{E}$ a function which attaches to every
$f\in\Fl(\cal{E})$ an inverse image functor for $f$ in $\cal{F}$,
namely~$f^*$. The cleavage is said to be \emph{normalized}
if $f=\id_S$ implies $f^*=\id_{\cal{F}_S}$. One calls
\emph{cloven category} (\resp \emph{normalized cloven
category})
a category $\cal{F}$ over~$\cal{E}$ equipped with a cleavage
(\resp with a normalized cleavage).
\end{definition}

It is clear that $\cal{F}$ admits a cleavage if and only if
$\cal{F}$ is prefibered over~$\cal{E}$, and then $\cal{F}$
admits a normalized cleavage. The set of cleavages on~$\cal{F}$
is in bijective correspondence with the set of subsets $K$ of
$\Fl(\cal{F})$ satisfying the following conditions:
\begin{enumerate}
\item[a)] The $\alpha\in K$ are cartesian morphisms.
\item[b)] For every morphism $f\colon T\to S$ in $\cal{E}$ and every
$\xi\in\Ob(\cal{F}_S)$, there exists a unique $f$-morphism in $K$, with
target $\xi$.
\end{enumerate}

For the cleavage defined by $K$ to be normalized, it is necessary and
sufficient that $K$ satisfy in addition the condition
\begin{enumerate}
\item[c)] The identity morphisms in $\cal{F}$ belong
to~$K$.
\end{enumerate}

% original p. 171
The morphisms that are elements of~$K$ may be called the
\emph{``transport morphisms''}
for the cleavage under consideration.

The notion of isomorphism of cloven categories over~$\cal{E}$
is clear. More generally, one may define the
morphisms of cloven $\cal{E}$-categories as the functors
of $\cal{E}$-categories $\cal{F}\to \cal{G}$ which send
transport morphisms to transport morphisms. (These are in
particular cartesian functors). In this way the
cloven categories over~$\cal{E}$ are the objects of a
category, the \emph{category of cloven categories
over} $\cal{E}$. The reader will spell out the existence of products,
related to the fact that if a category over~$\cal{E}$ is a product of
categories $\cal{F}_i$ over~$\cal{E}$ each equipped with a cleavage,
then $\cal{F}$ is equipped with a corresponding natural cleavage. One
also leaves to the reader to spell out the notion of change of base
in cloven categories.

We shall denote by $\alpha_f(\xi)$ the canonical morphism
$$
\alpha_f(\xi)\colon f^*(\xi)\to\xi.
$$
It is, as has been said, functorial in $\xi$, \ie one has a functorial
homomorphism
$$
\alpha_f\colon i_Tf^*\to i_S,
$$
where for every $S\in\Ob(\cal{E})$, $i_S$ denotes the inclusion
functor
$$
i_S\colon\cal{F}_S\to \cal{F}
$$

Consider now morphisms
$$
f\colon T\to S\quad\text{ and }\quad g\colon U\to T
$$
in $\cal{E}$, and let $\xi\in\Ob(\cal{F}_S)$, there then exists a
unique $U$-morphism
$$
c_{f,g}(\xi)\colon g^*f^*(\xi)\to (fg)^*(\xi)
$$
making
% original p. 172
the diagram commutative
$$
\xymatrix{
f^*(\xi) \ar[d]_{\alpha_f(\xi)} & & \ar[ll]_{\alpha_g(f^*(\xi))}
g^*(f^*(\xi)) \ar[d]^{c_{f,g}(\xi)} \\
\xi & & \ar[ll]^{\alpha_{fg}(\xi)} (fg)^*(\xi) }
$$
(by virtue of the definition of~$(fg)^*(\xi)$). For variable $\xi$,
this homomorphism is functorial, {i.e.}: one has a homomorphism
$$
c_{f,g}\colon g^*f^*\to (fg)^*
$$
of functors $\cal{F}_S\to\cal{F}_U$. Let us note at once:
\begin{proposition}
\label{VI.7.2} For the cloven category $\cal{F}$
over~$\cal{E}$ to be fibered, it is necessary and sufficient that the $c_{f,g}$
be isomorphisms.
\end{proposition}

One concludes from this, taking for $f$ an isomorphism, for $g$ its inverse,
and considering the isomorphisms $c_{f,g}$ and $c_{g,f}$:
\begin{corollary}
\label{VI.7.3} If $\cal{F}$ is a \emph{fibered}
cloven category over~$\cal{E}$, then for every isomorphism $f\colon T\to
S$ in $\cal{E}$, $f^*$ is an equivalence of categories
$\cal{F}_S\to \cal{F}_T$.
\end{corollary}

\begin{proposition}
\label{VI.7.4} Let $\cal{F}$ be a cloven category over~$\cal{E}$.
One has

$\qquad\qquad A)\qquad
\left\{\begin{array}{ccc}c_{f,\id_T}(\xi)&=&\alpha_{\id_T}(f^*(\xi))\\
c_{\id_S,f}(\xi)&=&f^*(\alpha_{\id_S}(\xi))
\end{array}\right.$

\smallskip
$\qquad\qquad B)\qquad c_{f,gh}(\xi)\cdot
c_{g,h}(f^*(\xi))=c_{fg,h}(\xi)\cdot h^*(c_{f,g}(\xi))$

\noindent(In
% original p. 173
these formulas, $f,g,h$ denote morphisms
$$
V\to U\to T\to S
$$
and $\xi$ an object of~$\cal{F}_S$).
\end{proposition}

The first and second relations, in the case of a normalized
cleavage, take the simpler form
\begin{equation*}
\label{eq:VI.7.A'}\quad A')\qquad {c_{f,\id_T}=\id_{f^*},\quad
c_{\id_S,f}=\id_{f^*}.\qquad\qquad\qquad\qquad\qquad\qquad}
\end{equation*}

As for the third, it is visualized by the
commutativity of the diagram
\begin{equation*}
\tag{$D$}
\label{eq:VI.7.D}
\begin{split}
\xymatrix@C=2.2cm{
h^*g^*f^*(\xi) \ar[r]^-{c_{g,h}(f^*(\xi))} \ar[d]_{h^*(c_{f,g}(\xi))}
& (gh)^*(f^*(\xi)) \ar[d]^{c_{f,gh}(\xi)} \\
h^*(fg)^*(\xi) \ar[r]_-{c_{fg,h}(\xi)} & (fgh)^*(\xi) }
\end{split}
\end{equation*}
In the case of fibered
categories
(where the $c_{f,g}$ are isomorphisms), this commutativity
can be expressed intuitively by the fact that \emph{the successive
use of isomorphisms of the form $c_{f,g}$ does not lead
to ``contradictory identifications''.} One can equally
write this formula without argument $\xi$, by using the
convolution product of homomorphisms of functors:
$$
c_{fg,h}\circ (h^**c_{f,g})=c_{f,gh}\circ(c_{g,h}*f^*).
$$

The proof of the first two formulas~\ref{VI.7.4} is
trivial; let us sketch that of the third. For this,
consider, in addition to the square $(D)$, the square
of homomorphisms:
\begin{equation*}
\tag{$D'$}
\label{eq:VI.7.D'}
\begin{split}
\xymatrix@C=2cm{
g^*f^*(\xi) \ar[r]^-{\alpha_g(f^*(\xi))}
\ar[d]_{c_{f,g}(\xi)} & f^*(\xi) \ar[d]^{\alpha_f(\xi)} \\
(fg)^*(\xi) \ar[r]_-{\alpha_{fg}(\xi)} & \xi }
\end{split}
\end{equation*}
which
% original p. 174
is commutative by definition of
$c_{f,g}(\xi)$. Consider the diagram obtained by joining the
vertices of~$(D)$ to the corresponding vertices of~$(D')$ by the
homomorphisms of the form $\alpha$:
$$
\begin{array}{ll}
\alpha_h(g^*f^*(\xi)),&\alpha_{gh}(f^*(\xi))\\
\alpha_h((fg)^*(\xi)),&\alpha_{fgh}(\xi).
\end{array}
$$
The four lateral faces of the cube thus obtained are likewise
commutative: for the left-hand one, this comes from the fact that the
left-hand column of~$(D)$ is deduced from the left-hand column of~$(D')$
by applying~$h$, and that $\alpha_h$ is a functorial
homomorphism; for the other three, it is nothing other than the
definition of the operations $c$ of the three remaining sides
of~$(D)$. Thus the five faces of the cube other than the
top face are commutative. It follows that the two
$(fgh)$-morphisms $h^*g^*f^*(\xi)\to(fgh)^*(\xi)$ defined by
$(D)$ have the same composite with $\alpha_{fgh}(\xi)\colon
(fgh)^*(\xi)\to\xi$, hence they are equal by definition of
$(fgh)^*$.

Let us restrict ourselves in what follows to \emph{normalized}
cloven categories. Such a category gives rise to the
following objects:
\begin{enumerate}
\item[a)] A map $S\mto \cal{F}_S$ from~$\Ob(\cal{E})$ to
$\Cat$.
\item[b)] A map $f\mto f^*$, associating to every
$f\in\Fl(\cal{E})$, with source $T$ and target $S$,
a functor $f^*\colon \cal{F}_S\to\cal{F}_T$.
\item[c)] A map $(f,g)\mto c_{f,g}$, associating to every
pair of arrows $(f,g)$ of~$\cal{E}$, a functorial
homomorphism $c_{f,g}\colon g^*f^*\to (fg)^*$.
\end{enumerate}

Moreover these data satisfy the conditions expressed in
the formulas $A')$ and~$B)$ given above. (N.B. If one had not
restricted oneself to the case of a normalized cleavage, it would have
been necessary to introduce an additional object, namely a function
$S\mto \alpha_S$ which associates to every object $S$ of~$\cal{E}$ a
functorial homomorphism $\alpha_S\colon (\id_S)^*\to
\id_{\cal{F}_S}$; the condition $A')$ would then be replaced by the
condition $A)$).

% original p. 175
We shall now show how one can reconstruct (up to
unique isomorphism) the normalized cloven
category $\cal{F}$ over~$\cal{E}$ with the help of the preceding
objects.

\section[Cloven category defined by a pseudofunctor]{Cloven category defined by a pseudofunctor
$\cal{E}^\circ\to\Cat$}
\label{VI.8}

Let us call, for brevity, a \emph{pseudofunctor} from~$\cal{E}^\circ$
to $\Cat$ (one should say, a \emph{normalized} pseudofunctor),
a set of data a), b), c) as above, satisfying the
conditions $A')$ and $B)$. In the preceding no., we
associated, to a normalized cloven category
over~$\cal{E}$, a pseudofunctor $\cal{E}^\circ\to\Cat$; here we
shall indicate the inverse construction. We shall leave to the reader the
verification of most of the details, as well as of the fact that
these constructions are indeed ``inverse'' to one another. More
precisely, one should consider the
pseudofunctors $\cal{E}^\circ\to\Cat$ as the objects of a
new category, and show that our constructions
furnish equivalences, quasi-inverse to one another,
between the latter and the category of
cloven categories over~$\cal{E}$, defined in the preceding
no.

One sets
$$
\cal{F}_\circ=\underset{S\in\Ob(\cal{E})}\coprod
\Ob(\cal{F}(S)),
$$
the sum of the sets
$\Ob(\cal{F}(S))$
(N.B. we shall write here $\cal{F}(S)$ and not $\cal{F}_S$ for the value at
the object $S$ of~$\cal{E}$ of the given pseudofunctor, to avoid
confusions of notation in what follows). One thus has an obvious
map:
$$
p_\circ\colon\cal{F}_\circ\to\Ob\cal{E}.
$$
Let
$$
\overline{\xi}
=(S,\xi),\quad \overline{\eta}=(T,\eta)\qquad\qquad
\text{(with $\xi\in\Ob\cal{F}(S)$, $\eta\in\Ob\cal{F}(T)$)}
$$
be two elements of~$\cal{F}_\circ$, and let $f\in\Hom(T,S)$, one
will set
$$
h_f(\overline{\eta},\overline{\xi})=\Hom_{\cal{F}(T)}(\eta,f^*(\xi)).
$$
% original p. 176
If moreover one has a morphism $g\colon U\to T$ in $\cal{E}$, and a
$\zeta\in\Ob\cal{F}(U)$, one defines a map, denoted
$(u,v)\mto u\circ v$:
$$
h_f(\overline{\eta},\overline{\xi})\times
h_g(\overline{\zeta},\overline{\eta})\to
h_{fg}(\overline{\zeta},\overline{\xi}),
$$
\ie a map
$$
\Hom_{\cal{F}(T)}(\eta,f^*(\xi))\times\Hom_{\cal{F}(U)}(\zeta,g^*(\eta))
\to \Hom_{\cal{F}(U)}(\zeta,(fg)^*(\xi)),
$$
by the formula
$$
u\circ v=c_{f,g}(\xi)\cdot g^*(u)\cdot v,
$$
\ie $u\circ v$ is the composite of the sequence
$$
\zeta\lto{u}
g^*(\eta)\lto{g^*(u)}
g^*f^*(\xi)\lto{c_{f,g}(\xi)} (fg)^*(\xi).
$$
On the other hand one will set
$$
h(\overline{\eta},\overline{\xi})=\underset{f\in\Hom(T,S)}\coprod
h_f(\overline{\eta},\overline{\xi}),
$$
and the preceding pairings define
pairings
$$
h(\overline{\eta},\overline{\xi})\times
h(\overline{\zeta},\overline{\eta})\to
h(\overline{\zeta},\overline{\xi}),
$$
while the definition of the $h(\overline{\eta},\overline{\xi})$
implies an obvious map:
$$
p_{\overline{\eta},\overline{\xi}}\colon
h(\overline{\eta},\overline{\xi})\to \Hom(T,S).
$$
This being said, one verifies the following points:
\begin{enumerate}
\item[1)] The composition between elements of the
$h(\overline{\eta},\overline{\xi})$
is \emph{associative}.
\item[2)]
% original p. 177
For every $\overline{\xi}=(\xi,S)$ in $\cal{F}_\circ$,
consider the identity element
of
$$
h_{\id_S}(\overline{\xi},\overline{\xi})=
\Hom_{\cal{F}(S)}(\id_S^*(\xi),\xi)=\Hom_{\cal{F}(S)}(\xi,\xi),
$$
and its image in $h(\overline{\xi},\overline{\xi})$. This object is
a \emph{unit} on the left and on the right for the composition
between elements of the $h(\overline{\eta},\overline{\xi})$.

This already shows that \emph{one obtains a category}
$\cal{F}$, by setting
$$
\Ob\cal{F}=\cal{F}_\circ,\quad \Fl{\cal{F}}=
\underset{\overline{\xi},\overline{\eta}\in\cal{F}_\circ}
\coprod h(\overline{\eta},\overline{\xi}).
$$
(N.B. one cannot simply take for $\Fl\cal{F}$ the
\emph{union} of the sets $h(\overline{\eta},\overline{\xi})$,
for the latter are not necessarily disjoint). Moreover:
\item[3)] The maps $p_\circ\colon \Ob\cal{F}\to\Ob\cal{E}$ and
$p_1=(p_{\overline{\eta},\overline{\xi}})\colon\Fl\cal{F}\to\Fl{\cal{E}}$
define a \emph{functor} $p\colon\cal{F}\to\cal{E}$. In
this way, $\cal{F}$ becomes a category over~$\cal{E}$, and
moreover the obvious map
$h_f(\overline{\eta},\overline{\xi})\to\Hom(\overline{\eta},\overline{\xi})$
induces a \emph{bijection}
$$
h_f(\overline{\eta},\overline{\xi})\isomto
\Hom_f(\overline{\eta},\overline{\xi}).
$$
\item[4)] The obvious maps
$$ \Ob\cal{F}(S)\to\cal{F}_\circ=\Ob\cal{F},\qquad \Fl\cal{F}(S)\to
\Fl\cal{F},
$$
where the second is defined by the obvious
maps
$$
\Hom_{\cal{F}(S)}(\xi,\xi')=
h_{\id_S}(\overline{\xi},\overline{\xi}')\to
\Hom(\overline{\xi},\overline{\xi}')
$$
define an \emph{isomorphism}
$$
i_S \colon \cal{F}(S)\isomto\cal{F}_S.
$$
\item[5)] For every object $\overline{\xi}=(S,\xi)$ of~$\cal{F}$, and
every morphism $f\colon T\to S$ of~$\cal{E}$, consider
% original p. 178
the element $\overline{\eta}=(T,\eta)$ of~$\cal{F}_T$, with
$\eta=f^*(\xi)$, and the element $\alpha_f(\xi)$ of
$\Hom(\overline{\eta},\overline{\xi})$, image of~$\id_{f^*(\xi)}$ by
the morphism
$\Hom_{\cal{F}(T)}(f^*(\xi),f^*(\xi))=h_f(\overline{\eta},\overline{\xi})
\to \Hom_f(\overline{\eta},\overline{\xi})$. \emph{This element
is cartesian, and it is the identity in $\overline{\xi}$ if}
$f=\id_S$, in other words, the set of the $\alpha_f(\xi)$
defines a \emph{normalized cleavage of~$\cal{F}$ over}
$\cal{E}$. Moreover, by construction, one has commutativity in the
diagram of functors
$$
\xymatrix{
\cal{F}(S) \ar[r]^{f^*} \ar[d]_{i_S} & \cal{F}(T) \ar[d]^{i_T} \\
\cal{F}_S \ar[r]^{f^*_{\cal{F}}} & \cal{F}_T }
$$
where $f^*_{\cal{F}}$ is the inverse image functor by $f$, relative
to the cleavage under consideration on~$\cal{F}$. Finally:
\item[6)] the homomorphisms $c_{f,g}$ given with the
pseudofunctor are transformed, by the isomorphisms $i_S$, into
the functorial homomorphisms $c_{f,g}$ associated to the cleavage
of~$\cal{F}$.
\end{enumerate}

We confine ourselves to giving the verification of 1) (which is, if
possible, less trivial than the others). It suffices to prove
the associativity of the composition between the objects of sets of
the form $h_f(\overline{\eta},\overline{\xi})$. Consider therefore
in $\cal{E}$ morphisms
$$
S\lfrom{f} T\lfrom{g} U\lfrom{h} V
$$
and objects
$$
\xi,\ \eta,\ \zeta,\ \tau
$$
in $\cal{F}(S),\cal{F}(T),\cal{F}(U),\cal{F}(V)$, finally
elements
$$
\begin{array}{ccc}
u\in
h_f(\overline{\eta},\overline{\xi})&=&\Hom_{\cal{F}(T)}(\eta,f^*(\xi))\\
v\in
h_g(\overline{\zeta},\overline{\eta})&=&\Hom_{\cal{F}(U)}(\zeta,g^*(\eta))\\
w\in
h_h(\overline{\tau},\overline{\zeta})&=&\Hom_{\cal{F}(V)}(\tau,h^*(\zeta)).
\end{array}
$$
One
% original p. 179
wants to prove the formula
$$
(u\circ v)\circ w = u\circ (v\circ w),
$$
which is an equality in
$\Hom_{\cal{F}(V)}(\tau,(fgh)^*(\xi))$. By virtue of the definitions
the two sides of this equality are obtained by composition
along the upper and lower contours of the diagram
below:
$$
\xymatrix{
\tau \ar[r]^w \ar[drr]_{v\circ w} & h^*(\zeta)
\ar[r]^{h^*(v)} \ar@<-2pt>
`u[rrrrr] `[rrrrr]^{h^*(u\circ v)} [rrrrr] & h^*g^*(\eta)
\ar[rr]^{h^*g^*(u)}
\ar[d]_{c_{g,h}(\eta)} & & h^*g^*f^*(\xi) \ar[rr]^{h^*(c_{f,g}(\xi))}
\ar[d]^{c_{g,h}(f^*(\xi))} & & h^*(fg)^*(\xi) \ar[d]^{c_{fg,h}(\xi)} \\
& & (gh)^*(\eta) \ar[rr]^{(gh)^*(u)} & & (gh)^*f^*(\xi)
\ar[rr]^{c_{f,gh}(\xi)} & & (fgh)^*(\xi) }
$$
Now the middle square is commutative because $c_{g,h}$ is a
functorial homomorphism, and the right-hand square is commutative by
virtue of the condition~$B)$ for a pseudofunctor. Whence the
announced result.

Of course, it remains to specify, when the pseudofunctor
under consideration already comes from a normalized cloven
category $\cal{F}'$ over~$\cal{E}$, how one obtains a
natural isomorphism between $\cal{F}'$ and $\cal{F}$. We leave the
details to the reader.

We also leave to the reader to interpret, in terms of
pseudofunctors, the notion of inverse image of a cloven
category $\cal{F}$ over~$\cal{E}$ by a change of base functor
$\cal{E}' \to\cal{E}$.

\section[Example]{Example: cloven category defined by a functor
$\cal{E}^\circ\to\Cat$; split categories over~$\cal{E}$}
\label{VI.9}
Suppose one has a functor
$$
\varphi\colon\cal{E}^\circ\to\Cat,
$$
it then defines a pseudofunctor by setting
$$
\cal{F}(S)=\varphi(S),
\quad f^*=\varphi(f),\quad c_{f,g}=\id_{(fg)^*}
$$
% original p. 180
Thus the construction of the preceding no. gives us a
category $\cal{F}$ cloven over~$\cal{E}$, said to be associated to the
functor $\varphi$. For a cloven category over~$\cal{E}$
to be isomorphic to a cloven category defined by a
functor $\varphi:\cal{E}^\circ\to\Cat$, it is necessary and sufficient,
manifestly, that it satisfy the conditions:
$$
(fg)^*=g^*f^*,\quad c_{f,g}=\id_{(fg)^*}\quoi.
$$
In terms of the set $K$ of transport morphisms, this also simply
means that \emph{the composite of two transport
morphisms is a transport morphism}. A cleavage of a
category $\cal{F}$ over~$\cal{E}$ satisfying the preceding
condition is called a \emph{splitting}
of~$\cal{F}$ over~$\cal{E}$, and a category $\cal{F}$
over~$\cal{E}$ equipped with a splitting is called a
\emph{split category over~$\cal{E}$}.
It is thus a particular case of the notion of
cloven category. The category of split categories
over~$\cal{E}$ is thus equivalent to
$\SheafHom(\cal{E}^\circ,\Cat)$. Note that a
split category over~$\cal{E}$ is a fortiori a cloven category
over~$\cal{E}$.

If $\cal{F}$ is a fibered category over~$\cal{E}$, there
does not always exist a splitting on~$\cal{F}$. Suppose for example
that $\Ob\cal{E}$ and $\Ob\cal{F}$ are reduced to a
single element, and that the set of endomorphisms of the said element is a
group $E$ \resp $F$, so that the projection functor $p$ is
given by a homomorphism of groups $p\colon F\to E$, surjective
since $p$ is fibrant. One then verifies at once that
the set of cleavages of~$\cal{F}$ over~$\cal{E}$ is in
bijective correspondence with the set of maps $s\colon
E\to F$ such that $ps=\id_E$ (\ie the set of ``systems of
representatives'' for the classes mod the subgroup $G$ kernel of
the surjective homomorphism $p\colon F\to E$). A cleavage is a
splitting if and only if $s$ is a homomorphism of groups. To say
that there exists a splitting therefore means that the group extension $F$
of~$E$ by $G$ is trivial, which is expressed, when $G$ is
commutative, by the vanishing of a certain cohomology class
in $\H^2(E,G)$ (where $G$ is regarded as a group
on which $E$ acts).

Suppose however that $\cal{F}$ is a fibered category
over~$\cal{E}$ such that the
% original p. 181
$\cal{F}_S$ are \emph{rigid} categories, \ie the group
of automorphisms of every object of~$\cal{F}_S$ is reduced to
the identity. It is then easy to prove that $\cal{F}$ admits a
splitting over~$\cal{E}$. Indeed, one first observes that the question
of existence of a splitting is not modified if one replaces
$\cal{F}$ by an $\cal{E}$-equivalent category, which in the present
instance brings us back to the case where the $\cal{F}_S$ are
rigid \emph{and reduced} categories (\ie two isomorphic
objects in $\cal{F}_S$ are identical). But if $G$ is a
rigid and reduced category, every isomorphism of two
functors $H\to G$ (where $H$ is an arbitrary category) is
an identity. It follows that if $\cal{F}$ is a
fibered category over~$\cal{E}$, such that the fiber-categories are
rigid and reduced, then there exists a \emph{unique} cleavage of
$\cal{F}$ over~$\cal{E}$, which is necessarily a splitting. Thus
$\cal{F}$ is isomorphic to the category defined by a
functor $\varphi\colon \cal{E}^\circ\to \Cat$, such that the
$\varphi(S)$ are rigid and discrete categories, and the
functor $\varphi$ is defined up to isomorphism.

% ---- en-chunk4.tex ----

\section{Cofibered categories, bifibered categories}
\label{VI.10}
Let us consider a category $\cal{F}$ over~$\cal{E}$,
with the projection functor
$$
p\colon \cal{F}\to \cal{E},
$$
it defines a category $\cal{F}^\circ$ over
$\cal{E}^\circ$, by the projection functor
$$
p^\circ\colon \cal{F}^\circ\to \cal{E}^\circ.
$$
A morphism $\alpha\colon \eta\to\xi$ in $\cal{F}$ is said to be
\emph{cocartesian} if it is a cartesian morphism for
$\cal{F}^\circ$ over~$\cal{E}^\circ$. Making this explicit, one sees that this
means that for every object $\xi'$ of~$\cal{F}_S$, the map
$u\mto u\circ \alpha$
$$
\Hom_S(\xi,\xi')\to\Hom_f(\eta,\xi')
$$
is bijective. One then also says that $(\xi,\alpha)$ is a
\emph{direct image}
of~$\eta$ by $f$, in the category $\cal{F}$ over~$\cal{E}$. If
it exists for every $\eta$ in $\cal{F}_T$, one says that the
direct image functor by $f$ exists, and one denotes this functor
% original p. 182
$f_*^{\cal{F}}$ or~$f_*$,
once it has been chosen. It is thus defined by
an isomorphism of bifunctors on~$\cal{F}_T^\circ\times\cal{F}_S$:
$$
\Hom_S(f_*(\eta),\xi)\isomto \Hom_f(\eta,\xi).
$$
Thus if $f_*$ exists, for $f^*$ to exist, it is necessary and sufficient that
$f_*$
admit an adjoint functor, \ie that there exist a functor
$f^*\colon \cal{F}_S\to\cal{F}_T$ and an isomorphism of bifunctors
$$
\Hom_S(f_*(\eta),\xi)\isomto \Hom_T(\eta,f^*(\xi)).
$$
Let $g\colon U\to T$ be another morphism in $\cal{E}$, and suppose
that the inverse images and direct images by $f,g$ and $fg$
exist. Let us consider then the functorial homomorphisms
$$
\begin{array}{cccc}
c^{f,g}\colon&f_*g_*&\dpl\from&(fg)_*\\
c_{f,g}\colon&g^*f^*&\dpl\to&(fg)^*.
\end{array}
$$
One observes that if one regards $f_*g_*$ and $g^*f^*$ as a
pair of adjoint functors, as well as $(fg)_*$ and $(fg)^*$, the two
preceding homomorphisms are adjoint to one another. Thus
one is an isomorphism if and only if the other is. In
particular:
\begin{proposition}
\label{VI.10.1} Suppose that the category $\cal{F}$ over~$\cal{E}$
is prefibered and coprefibered. For it to be
fibered, it is necessary and sufficient that it be
cofibered.
\end{proposition}

Of course, one says that $\cal{F}$ is coprefibered \resp cofibered
over~$\cal{E}$, if $\cal{F}^\circ$ is prefibered \resp fibered over~$\cal{E}$. We shall say that $\cal{F}$ is bifibered
over~$\cal{E}$,
if it is at once fibered and cofibered over~$\cal{E}$.

\section{Various examples}
\label{VI.11}

\begin{enumerate}

\item[a)] \textbf{Categories of arrows of}~$\cal{E}$. Let
$\cal{E}$ be a category. Let us denote by $\mathbf{\Delta^1}$ the
category associated to the totally ordered set
with two elements $[0,1]$;
% original p. 183
it thus has two objects 0 and 1, and besides the two identity
morphisms an arrow $(0,1)$ with source 0 and target 1. Let
$$
\CatFl(\cal{E})=\SheafHom(\mathbf{\Delta^1},\cal{E})
$$
one calls it the \emph{category of arrows of} $\cal{E}$.
The object 1 of~$\mathbf{\Delta^1}$ defines a canonical functor,
called the \emph{target functor}
$$
\CatFl(\cal{E})\to \cal{E}
$$
(the functor defined by the object 0 of~$\mathbf{\Delta^1}$ is
called the \emph{source functor}). For every object $S$ of~$\cal{E}$,
the fiber-category $\CatFl(\cal{E})_S$ is canonically isomorphic
to the category $\cal{E}_{/S}$ of objects of~$\cal{E}$
over~$S$.

Let us consider a morphism $f\colon T\to S$ in $\cal{E}$, then there
corresponds to it a canonical functor
$$
f_*\colon \cal{E}_{/T}=\cal{F}_T\to \cal{E}_{/S}=\cal{F}_S
$$
and a functorial isomorphism
$$
\Hom_S(f_*(\eta),\xi)\isomto\Hom_f(\eta,\xi)
$$
which thus makes of~$f_*$ a direct image functor for $f$ in
$\cal{F}$. One has moreover here
$$
(\id_S)_*=\id_{\cal{F}_S},\quad (fg)_*=f_*g_*,\quad
c^{f,g}=\id_{(fg)},
$$
\ie $\cal{F}$ is equipped with a co-splitting over~$\cal{E}$. A fortiori,
$\cal{F}$ is cofibered over~$\cal{E}$. Let us now note that
the set of morphisms in $\cal{F}$ is in bijective
correspondence with the set of commutative square diagrams in
$\cal{E}$.
$$
\xymatrix{
X\ar[d]_-u&\ar[l]_-{f'}Y\ar[d]^-v\\
S&\ar[l]_-{f}T
}
$$
By
% original p. 184
definition, the morphism in question is cartesian if the
square is cartesian in $\cal{E}$, \ie if it makes of~$Y$ a
fibered product of~$X$ and $T$ over~$S$. The inverse image functor
$f^*$
thus exists if and only if for every object $X$ over~$S$, the fibered
product $X\times_ST$ exists. It follows from~\ref{VI.10.1} that if
the product of two objects over a third always exists in~$\cal{E}$, \ie if $\cal{F}$ is prefibered over~$\cal{E}$,
then $\cal{F}$ is even fibered over~$\cal{E}$.

\medskip
\item[b)] \textbf{Category of presheaves or sheaves on
variable spaces}

Let $\cal{E}=\mathbf{Top}$ be the category of topological spaces.
If $T$ is a topological space, we shall denote by $\cal{U}(T)$ the
category of open sets of~$T$, where the morphisms are the
inclusion maps. If $\cal{C}$ is a category, a
functor $\cal{U}(T)^\circ\to \cal{C}$ is called a
\emph{presheaf} on~$T$ with values in $\cal{C}$, and a
\emph{sheaf} if it satisfies a left exactness condition
which we do not repeat here. The \emph{category
$\cal{P}(T)$ of presheaves on~$T$ with values in
$\cal{C}$} is by definition the category
$\SheafHom(\cal{U}(T)^\circ,\cal{C})$, and the category
$\cal{F}(T)$ of sheaves on~$T$ with values in $\cal{C}$ is the
full subcategory whose objects are the objects of
$\SheafHom(\cal{U}(T)^\circ,\cal{C})$ which are sheaves. If
$f\colon T\to S$ is a morphism in $\cal{E}$, \ie a
continuous map of topological spaces, there corresponds to it by
the increasing map $U\mto f^{-1}(U)$ a functor
$\cal{U}(S)\to \cal{U}(T)$, whence a functor
$$
f_*\colon\SheafHom(\cal{U}(T)^\circ,\cal{C})\to
\SheafHom(\cal{U}(S)^\circ,\cal{C})
$$
called the \emph{direct image functor of presheaves by}
$f$. One sees immediately that the direct image of a sheaf is a
sheaf, thus the functor
$f_*\colon \cal{P}(T)\to\cal{P}(S)$
induces a functor, likewise denoted
$f_*\colon\cal{F}(T)\to\cal{F}(S)$.
One moreover verifies trivially (by the associativity of the
composition of functors) that one has, for a second continuous
map $g\colon U\to T$, the identity
$$
(gf)_* =g_*f_*,\qquad
\text{likewise }\ (\id_S)_*=\id_{\cal{P}(S)}.
$$
In this way, one has obtained a functor
$$
S\mto \cal{P}(S)
$$
\resp $$
S\mto \cal{F}(S)
$$
from
% original p. 185
$\cal{E}$ to~$\Cat$. In fact, we are interested in the
corresponding functor
$$
S\mto \cal{P}(S)^\circ,\qquad\text{\resp~}S\mto
\cal{F}(S)^\circ.
$$
It defines a cofibered category, and even a
co-split one, over the category of topological spaces, which one
calls the \emph{cofibered category of presheaves}
(\resp \emph{sheaves}) \emph{with values in} $\cal{C}$
(understood: on variable spaces). Making explicit the construction
of~no.~\ref{VI.8}, one sees that a morphism from a presheaf $B$ on~$T$
to a presheaf $A$ on~$S$ is a pair $(f,u)$ consisting
of a continuous map from~$T$ to $S$, and of a morphism $u\colon
A\to f_*(B)$ in the category $\cal{P}(S)$. This description
is equally valid for morphisms of sheaves, $\cal{F}$
being a full subcategory of~$\cal{P}$.

In the most important cases, the category $\cal{P}$ and the
category $\cal{F}$ over~$\cal{E}$ are also
fibered categories, \ie for every continuous map, the
direct image functors $\cal{P}(T)\to \cal{P}(S)$ and
$\cal{F}(T)\to\cal{F}(S)$ have an adjoint functor, which is then
denoted $f^*$ and called the inverse image functor of
presheaves \resp the inverse image functor of sheaves, by
the continuous map $f$. This functor exists for example if
$\cal{C}=\Ens$. One can show that the functor $f^*\colon
\cal{P}(S)\to \cal{P}(T)$ exists whenever in $\cal{C}$ the
inductive limits (relative to diagrams in the Universe
under consideration) exist. The question is less easy for
$\cal{F}$; indeed, one notes that (even in the case
$\cal{C}=\Ens$) the inverse image of a presheaf which is a
sheaf is in general not a sheaf, in other words the
inverse image functor of sheaves is not isomorphic to the functor
induced by the inverse image functor of presheaves (despite
the common notation $f^*$). Thus, $\cal{F}$ is a
cofibered subcategory of~$\cal{P}$, but not a
fibered subcategory, \ie \emph{the inclusion functor
$\cal{F}\to\cal{P}$ is not fibrant}.

The cofibered category $\cal{P}$ can be deduced from a more
general cofibered (or rather fibered) category,
obtained as follows. For every category $\cal{U}$ (in
the fixed Universe), one sets
$$
\cal{P}(\cal{U})=\SheafHom(\cal{U},\cal{C})
$$
and
% original p. 186
one notes that $\cal{U}\mto\cal{P}(\cal{U})$ is in a
natural way a contravariant functor in~$\cal{U}$, from the category
$\Cat$ to $\Cat$. It thus defines a split category
over $\cal{E}=\Cat$, which we shall denote $\Cat_{/\!/\cal{C}}$. The
objects of this category are the pairs $(\cal{U},p)$ of a
category $\cal{U}$ and a functor $p\colon\cal{U}\to\cal{C}$,
and a morphism from~$(\cal{U},p)$ to $(\cal{V},q)$ is
essentially a pair $(f,u)$, where $f$ is a functor
$\cal{U}\to\cal{V}$ and $u$ a homomorphism of functors $u\colon p\to
qf$. We leave to the reader the task of making explicit the composition of
morphisms in $\Cat_{/\!/\cal{C}}$. The projection functor
$$
\cal{F}=\Cat_{/\!/\cal{C}}\to\cal{E}=\Cat
$$
associates to the pair $(\cal{U},p)$ the object $\cal{U}$; the
fiber-category at $\cal{U}$ is the category
$\SheafHom(\cal{U},\cal{C})$ (up to isomorphism). When
inductive limits exist in $\cal{C}$, one shows easily
that the fibered category $\Cat_{/\!/\cal{C}}$ over~$\Cat$ is
likewise cofibered over~$\Cat$, \ie one can define the
notion of \emph{direct image of a functor} $p\colon\cal{U}\to\cal{C}$
by a functor $f\colon\cal{U}\to\cal{V}$. The category of
presheaves is deduced from the preceding fibered category
by the change of base
$$
\mathbf{Top}^\circ\to\Cat
$$
(functor $S\mto\cal{U}(S)$ defined above), which gives a
fibered category over~$\mathbf{Top}^\circ$, and passing to
the opposite category, one obtains the
cofibered category $\cal{P}$ of presheaves over
$\mathbf{Top}$. The notion of inverse image of a functor corresponds
to that of direct image of a presheaf, the notion of direct
image of a functor to that of inverse image of a
presheaf.

\medskip
\item[c)] \textbf{Objects with operators over an object
with operators}

Let $\cal{F}$ be a category over~$\cal{E}$, and let $S$ be an object of
$\cal{E}$ on which a group $G$ operates, on the left to fix
ideas. This object with operators can be interpreted as
corresponding to a functor $\lambda\colon\cal{E}'\to \cal{E}$ from
the category (with a single object, having $G$ as group
of endomorphisms) $\cal{E}'$ defined by $G$, into the
category $\cal{E}$, and thus there is defined by change of base
a category $\cal{F}'$ over~$\cal{E}'$, which is
fibered \resp cofibered when $\cal{F}$ is so
over~$\cal{E}$. A
% original p. 187
section of~$\cal{E}'$ over~$\cal{F}'$ (necessarily
cartesian, since $\cal{E}'$ is a groupoid, and every
isomorphism in $\cal{F}'$ is cartesian by virtue
of~\ref{VI.6.12}), can also be interpreted as an
$\cal{E}$-functor $\cal{E}'\to\cal{F}$ over~$\lambda$, or
also as an object with operators $\xi$ in $\cal{F}$
``over'' the object with operators $S$.

\medskip
\item[d)] \textbf{Pairs of quasi-inverse adjoint functors;
autodualities}

When the base-category $\cal{E}$ is reduced to two
objects $a,b$ and, besides the identity arrows, to two
isomorphisms $f\colon a\to b$ and $g\colon b\to a$ inverse to one
another (\ie $\cal{E}$ is a connected rigid groupoid with two
objects), a normalized cloven category over~$\cal{E}$ is
essentially the same thing as the system formed by
two categories $\cal{F}_a$ and $\cal{F}_b$ and a \emph{pair of
adjoint functors} $G\colon \cal{F}_a\to\cal{F}_b$ and $F\colon
\cal{F}_b\to\cal{F}_a$, which are equivalences of
categories (hence quasi-inverse to one another). One will take for
$\cal{F}_a$ and $\cal{F}_b$ the fiber-categories of~$\cal{F}$,
for $F$ and $G$ the functors $f^*$ and $g^*$, and the two
isomorphisms
$$
u\colon FG\isomto\id_{\cal{F}_a}\qquad v\colon
GF\isomto\id_{\cal{F}_b}
$$
are $c_{g,f}$ and $c_{f,g}$. The two usual
compatibility conditions between $u$ and $v$ are none other than the condition
\ref{VI.7.4}~$B)$ for the composites $fgf$ and $gfg$. It is easy
to show that these conditions suffice to imply that one indeed has
a pseudofunctor $\cal{E}^\circ\to\Cat$.

An interesting case is that where one has
$$
\cal{F}_b=\cal{F}_a^\circ,\quad G=F^\circ,\quad v=u^\circ.
$$
One calls \emph{autoduality}
in a category $\cal{C}$, the
data of a functor $D\colon \cal{C}\to\cal{C}^\circ$, and of an
isomorphism $u\colon DD^\circ\isomto\id_{\cal{C}}$, such that $u$ and
the isomorphism $u^\circ\colon D^\circ D\isomto\id_{\cal{C}^\circ}$
make of~$(D,D^\circ)$ a pair of adjoint functors,
(necessarily quasi-inverse to one another). This condition
is written:
$$
D(u(x))=u(D(x))\qquad \text{for every $x\in\Ob(\cal{C})$.}
$$
\item[e)]
% original p. 188
\textbf{Categories over a discrete category $\cal{E}$.} One says that $\cal{E}$ is a
\emph{discrete category} if every arrow therein is an
identity arrow, so that $\cal{E}$ is defined up to
unique isomorphism by the knowledge of the set
$I=\Ob(\cal{E})$. The data of a category $\cal{F}$
over~$\cal{E}$ is thus equivalent (up to unique isomorphism)
to the data of a family of categories
$\cal{F}_i$ ($i\in I$), the fiber-categories. Every category
$\cal{F}$ over~$\cal{E}$ is fibered, every $\cal{E}$-functor
$\cal{F}\to\cal{G}$ is cartesian, one has a canonical isomorphism
$$ \SheafHom_{\cal{E}/-}(\cal{F},\cal{G}) \isomto
\prod_{i}\SheafHom(\cal{F}_i,\cal{G}_i).
$$
In particular, one obtains
$$
\mathbf{\Gamma}(\cal{F}/\cal{E})=\varprojLim\cal{F}/\cal{E}\isomto\prod_{i} \cal{E}_i.
$$
\item[f)] \textbf{Suppose that $\cal{E}$ has exactly two objects
$S$ and $T$,} and besides the \textbf{identity morphisms, a
morphism }$f\colon T\to S$. Then a category $\cal{F}$
over~$\cal{E}$ is defined, up to unique $\cal{E}$-isomorphism,
by the data of two categories $\cal{F}_S$
and $\cal{F}_T$ and of a bifunctor $H(\eta,\xi)$ on~$\cal{F}_T^\circ\times\cal{F}_S$, with values in $\Ens$. Indeed,
if~$\cal{F}$ is a category over~$\cal{E}$, one associates to it
the two fiber-categories $\cal{F}_S$ and~$\cal{F}_T$ and
the bifunctor $H(\eta,\xi)=\Hom_f(\eta,\xi)$. One leaves to the reader the
task of making explicit the construction in the opposite direction. For the
category under consideration to be fibered (or prefibered,
which amounts to the same) it is necessary and sufficient that the functor $H$ be
representable with respect to the argument $\xi$. For it
to be cofibered, it is necessary and sufficient that $H$ be
representable with respect to the argument~$\eta$.
\item[g)] Let $\cal{F}=\cal{C}\times\cal{E}$, considered
as a category over~$\cal{E}$ by means of
$\pr_2$. Then $\cal{F}$ is fibered and cofibered over~$\cal{E}$,
and is even equipped with a canonical splitting and co-splitting,
corresponding to the constant functor on~$\cal{E}$, \resp on~$\cal{E}^\circ$, with values in $\Cat$, of value $\cal{C}$. One has
$$
\mathbf{\Gamma}(\cal{F}/\cal{E})\simeq\SheafHom(\cal{E},\cal{C})
$$
and
% original p. 189
$\varprojLim\cal{F}/\cal{E}$ corresponds to the
full subcategory formed of the functors
$F\colon\cal{E\to\cal{C}}$ transforming arbitrary morphisms into
isomorphisms.
\end{enumerate}

\section{Functors on a cloven category}
\label{VI.12}

Let $\cal{F}$ be a normalized cloven category
over~$\cal{E}$. For every object $S$ of~$\cal{E}$ one denotes by
$$
i_S\colon \cal{F}_S\to\cal{F}
$$
the inclusion functor. One thus has a functorial homomorphism, for
every morphism $f\colon T\to S$ in~$\cal{E}$:
$$
\alpha_f\colon i_T f^*\to i_S,
$$
where $f^*$ is the change of base functor
$\cal{F}_S\to\cal{F}_T$ for $f$ defined by the cleavage. Let
now
$$
F\colon\cal{F}\to\cal{C}
$$
be a functor from~$\cal{F}$ to a category $\cal{C}$; set,
for every~$S\in\Ob(\cal{E})$,
$$
F_S=F\circ i_S\colon\cal{F}_S\to\cal{C}
$$
and for every~$f\colon T\to S$ in~$\cal{E}$,
$$
\varphi_f=F*\alpha_f\colon F_T f^*\to F_S
$$
One has thus, to every functor $F\colon \cal{F}\to\cal{C}$,
associated a family $(F_S)$ of functors $\cal{F}_S\to\cal{C}$, and
a family $(\varphi_f)$ of homomorphisms of functors $F_T f^*\to
F_S$. These families satisfy the following conditions:

a) $\varphi_{\id_S}=\id_{F_S}$.

b) For two morphisms $f\colon T\to S$ and $g\colon U\to T$
in~$\cal{E}$, one has commutativity in the square
of functorial homomorphisms:
$$
\xymatrix@C=2cm{ F_U g^* f^* \ar[r]^-{F_U*c_{f,g}} \ar[d]_{\varphi_g * f^*} &
F_U(fg)^* \ar[d]^{\varphi_{fg}} \\ F_T f^*\ar[r]^-{\varphi_f} & F_S.}
$$
% original p. 190
The first relation is trivial, and the second relation
is obtained by applying the functor~$F$ to the commutative diagram
$$
\xymatrix@C=2cm{ g^* f^* (\xi) \ar[r]^-{c_{f,g}(\xi)}
\ar[d]_{\alpha_g(f^*(\xi))} & (fg)^*(\xi) \ar[d]^{\alpha_{fg}(\xi)}\\
f^*(\xi) \ar[r]^-{\alpha_f(\xi)}& \xi }
$$
for a variable object $\xi$ in~$\cal{F}_S$.

If $G$ is a second functor $\cal{F}\to\cal{C}$, giving
rise to functors $G_S\colon\cal{F}_S\to\cal{C}$ and
functorial homomorphisms $\psi_f\colon G_Tf\to G_S$, and if $u\colon
F\to G$ is a functorial homomorphism, then there correspond to it
functorial homomorphisms $u*i_S$:
$$
u_S\colon F_S\to G_S
$$
and one observes immediately that for every morphism $f\colon T\to S$
in $\cal{E}$, one has commutativity in the squares
$$
\begin{array}{c}
\xymatrix@C=1.5cm{ F_Tf^* \ar[r]^-{\phi_f} \ar[d]_{u_T * f^*} & F_S
\ar[d]^{u_S} \\ G_T f^* \ar[r]^-{\psi_f}& G_S}
\end{array}
\leqno\quad \textup{c)}
$$

\begin{proposition}
\label{VI.12.1}
Let $\cal{H}(\cal{F},\cal{C})$ be the category whose objects are
the pairs of families $(F_S)$ ($S\in\Ob(\cal{F})$) of functors
$\cal{F}_S\to\cal{C}$, and of families $(\varphi_f)$
($f\in\Fl(\cal{F})$) of functorial homomorphisms $F_Tf^*\to F_S$,
satisfying the conditions \emph{a) and~b)}, and where the morphisms
are the families $(u_S)$ ($S\in\Ob(\cal{F})$) of homomorphisms
% original p. 191
$F_S\to G_S$, verifying the commutativity condition
\emph{c)} written above, (the composition of morphisms being
effected by the composition of the homomorphisms of functors
$\cal{F}_S\to\cal{C}$). Then the two laws made explicit above
define an \emph{isomorphism} $K$ of the category
$\SheafHom(\cal{F},\cal{C})$ with the category
$\cal{H}(\cal{F},\cal{C})$.
\end{proposition}

It is trivial that one has indeed here a \emph{functor} from the
first category to the second. This functor is fully
faithful, since for given $F,G$, $\Hom(F,G)\to\allowbreak\Hom(K(F),K(G))$
is trivially injective; to show that it is surjective, it suffices
to note that the commutativity condition c) expresses the
functoriality of the maps $u(\xi)=u_S(\xi)\colon
F(S)=F_S(\xi)\to G(\xi)=G_S(\xi)$ for the homomorphisms of the form
$\alpha_f(\xi)$ in $\cal{F}$, on the other hand one has functoriality
on each fiber-category \ie for the morphisms in $\cal{F}$
which are $T$-morphisms ($T\in\Ob(\cal{E})$), whence
functoriality for every morphism in $\cal{F}$,
since an
$f$-morphism (where $f\colon T\to S$ is a morphism in
$\cal{E}$
is in a unique way) a composite of a morphism
$\alpha_f(\xi)$ and a $T$-morphism. It remains therefore to prove that
the functor~$K$ is bijective on objects. The preceding
argument already shows that $K$ is injective on
objects; it remains to prove that it is surjective, \ie that if one starts
from a system $(F_S),(\varphi_f)$, satisfying a) and b); and if one
defines a map $\Ob\cal{F}\to\Ob\cal{C}$ by
$$
F(\xi)=F_S(\xi) \quad\text{for }
\xi\in\Ob\cal{F}_S\subset\Ob\cal{F}
$$
and a map $\Fl(\cal{F})\to\Fl(\cal{C})$ by
$$
F(\alpha_f(\xi)u')=\varphi_f(\xi)\,F_T(u')
$$
for every morphism $f\colon T\to S$ in $\cal{E}$, every object $\xi$
of~$\cal{F}_S$ and every $T$-morphism $u'$ with target $f^*(\xi)$, then one
obtains a \emph{functor} $F$ from~$\cal{F}$ to $\cal{C}$. Indeed,
the relation $F(\id_\xi)=\id_{F(\xi)}$ is trivial, it remains to
prove the multiplicativity $F(uv)=F(u)F(v)$ when one has an
$f$-morphism $u\colon\eta\to\xi$ and a $g$-morphism
$v\colon\zeta\to\nu$, with $f\colon T\to S$ and $g\colon U\to T$
morphisms of~$\cal{E}$. Setting $w=uv$, one will have
$$
u=\alpha_f(\xi)u'\quoi,\quad v=\alpha_g(\eta)v'\quoi,\quad w=\alpha_{fg}(\xi)w'
$$
% original p. 192
with
$$
w'=c_{f,g}(\xi)g^*(u')v'\quad \hbox{(\cf no.~\ref{VI.8})}.
$$

With these notations, one must prove the commutativity of the outer
contour of the diagram below:
$$
\xymatrix{ F_U(\zeta) \ar@{-<} `u[r] `[rrrrrr]^{F(w')} [rrrrrr]
\ar[rr]^{F_u(v')} \ar[drr]_{F(v)} & & F_Ug^*(\eta)
\ar[rr]^{F_ug^*(u')} \ar[d]^{\varphi_g(\eta)} & & F_Ug^*f^*(\xi)
\ar[rr]^{F_U(c_{f,g}(\xi))}\ar[d]^{\varphi_g(f^*(\xi))}& &
F_U(fg)^*(\xi) \ar[d]^{\varphi_{fg}(\xi)} \\ & & F_T(\eta) \ar@{-<}
`d[r] `[rrrr]_{F(u)} [rrrr] \ar[rr]_{F_T(u')} & & F_Tf^*(\xi)
\ar[rr]_{\varphi_f(\xi)} & & F_S(\xi) }
$$
Now the left triangle is commutative by definition of~$F(v)$, the
middle square is commutative since it is deduced from the homomorphism
$u'\colon \xi\to f(\eta)$ by the functorial homomorphism $\alpha_g$,
finally the right-hand square is commutative by virtue of condition
b). The desired conclusion follows.

Suppose now that $\cal{C}$ is likewise a
normalized cloven category over~$\cal{E}$, which we
shall henceforth call $\cal{G}$, and that we are interested
in the $\cal{E}$-functors from~$\cal{F}$ to $\cal{G}$. If $F$ is such a
functor, it induces functors
$$
F_S\colon \cal{F}_S\to\cal{G}_S
$$
for the fiber-categories. On the other hand, for every morphism
$f\colon T\to S$ in~$\cal{E}$, and every object $\xi$ in~$\cal{F}_S$,
the $f$-morphism $F(\alpha_f(\xi))$ factors in a unique way
through a $T$-morphism
$$
\varphi_f(\xi)\colon F_T(f^*_{\cal{F}}(\xi))\to
f^*_{\cal{G}}(F_S(\xi))
$$
(where the $\cal{F}$ or the $\cal{G}$ in the index indicates the
cloven category for which one takes the inverse image
functor), whence a functorial homomorphism of functors from
$\cal{F}_S$ to $\cal{G}_T$:
$$
\varphi_f\colon F_T f^*_{\cal{F}}\to f^*_{\cal{G}}F_S.
$$
% original p. 193
The two systems $(F_S)$ and $(\varphi_f)$ satisfy the
following conditions:

a') $\varphi_{\id_S}=\id_{F_S}$.

b') For two morphisms $f\colon T\to S$ and $g\colon U\to T$ in
$\cal{E}$, one has commutativity in the following diagram of functorial
homomorphisms:
$$
\xymatrix{ F_Ug_F^*f_F^* \ar[rr]^{F_U *c_{f,g}^\cal{F}}
\ar[d]_{\varphi_{g^*}f^*_{\cal{F}}} & & F_U(fg)^*_F
\ar[dd]^{\varphi_{fg}} \\ g^*_{\cal{G}}F_Tf^*_F
\ar[d]_{g^*_{\cal{G}}*\varphi_f} & & \\ g^*_{\cal{G}}f^*_{\cal{G}}F_S
\ar[rr]^{c_{f,g}^\cal{G}*F_S} & & (fg)^*_{\cal{G}}F_S. }
$$
We leave the verification to the reader, as well as
the statement and the proof of the analogue of
proposition~\ref{VI.12.1}, implying that one thus obtains a
bijective correspondence between the set of $\cal{E}$-functors from
$\cal{F}$ to $\cal{G}$, and the set of systems
$(F_S)$,$(\varphi_f)$ satisfying the conditions a') and b')
above. Of course, in this correspondence, the cartesian
functors are characterized by the property that
the homomorphisms $\varphi_f$ are isomorphisms.

\begin{remark}
Of course, it is most often advantageous to reason
directly on fibered categories without using
explicit cleavages, which in particular dispenses with having to appeal,
for the simple notion of~$\cal{E}$-functor or of cartesian $\cal{E}$-functor,
to a cumbersome interpretation as
above. It is to avoid unbearable heaviness, and to
obtain more intrinsic statements,
% original p. 194
that we have had to renounce starting as in \cite{VI.2} from the
notion of cloven category (called ``fibered
category'' in \loccit), which recedes to second rank in favor of
that of fibered category. It is moreover probable that,
contrary to the usage still predominant now,
tied to old habits of thought, it will eventually
prove more convenient in universal problems not to
put the emphasis on \emph{one} solution supposed chosen once
and for all, but to put all solutions on an equal
footing.
\end{remark}

\begin{thebibliography}{9}
\bibitem[1]{VI.1} A\ptbl \textsc{Grothendieck}, Sur quelques points d'alg\`ebre
homologique,
T\^ohoku Math.\ J. \textbf{9} (1957) p\ptbl 119--221.


\bibitem[2]{VI.2} A\ptbl \textsc{Grothendieck}, Technique de descente et
Th\'eor\`emes d'existence,~I, S\'eminaire Bourbaki~190,
December~1959.
\end{thebibliography}



\end{document}
